\IfFormatAtLeastTF{2025-06-01}{\DocumentMetadata{lang=en-US,pdfversion=2.0,tagging=on}}{}
\documentclass[11pt]{amsart}
\usepackage[T1]{fontenc}
\IfFileExists{lmodern.sty}{\usepackage{lmodern}}{\PassOptionsToPackage{expansion=false}{microtype}}
\usepackage{amsmath,amssymb,amsthm,mathtools}
\usepackage{etoolbox}
\usepackage[letterpaper,textwidth=5.6in,textheight=\textheight,hmarginratio=1:1]{geometry}
\usepackage[nopatch=footnote]{microtype}
\usepackage{xcolor}
\usepackage[colorlinks=true,linkcolor=blue,citecolor=blue,urlcolor=blue]{hyperref}
\newif\ifanonymous
\anonymousfalse
\ifanonymous
  \newcommand{\paperauthor}{Anonymous}
\else
  % Repository locator for the public (non-anonymous) build of
% non_mf_groups_exist.tex.  Kept in a separate file so that the anonymous
% build has an anonymous source archive and not merely an anonymous PDF:
% exclude this file from a double-blind submission and the manuscript still
% compiles, printing "supplied with this submission as supplementary
% material" in place of the link.
%
% Loaded only when \anonymousfalse is set in the preamble.
\newcommand{\coderepository}{https://github.com/SauersML/group-approximation}
\newcommand{\coderepositoryname}{github.com/SauersML/group-approximation}

  \newcommand{\paperauthor}{\paperauthorname}
\fi
\hypersetup{pdftitle={Non-MF Groups},pdfauthor={\paperauthor},pdfsubject={MF groups and operator-norm matrix approximation},pdfkeywords={MF groups, norm matrix coronas, MF radicals, property (T), asymptotic representations}}
\newcommand{\doi}[1]{\href{https://doi.org/#1}{\nolinkurl{doi:#1}}}
% Machine-readable source metadata attaching a printed statement to an
% accompanying Lean~4 declaration.  The macro is intentionally invisible in
% the PDF; the web edition and audit scripts consume its two arguments.
% Each \leanverified marker identifies a selected printed statement and a Lean
% declaration intended to prove it.  The markers do not assert line-by-line
% verification of the prose or complete sentence-census coverage of the
% manuscript.
\newcommand{\leanverified}[2]{\ignorespaces}
\allowdisplaybreaks
\raggedbottom
\emergencystretch=1.5em
\newtheorem{theorem}{Theorem}[section]
\newtheorem{proposition}[theorem]{Proposition}
\newtheorem{lemma}[theorem]{Lemma}
\newtheorem{corollary}[theorem]{Corollary}
\theoremstyle{plain}
\newtheorem{mainthm}{Theorem}
\renewcommand{\themainthm}{\Alph{mainthm}}
\newcommand{\F}{\mathbb F}
\newcommand{\C}{\mathbb C}
\newcommand{\GL}{\operatorname{GL}}
\newcommand{\EL}{\operatorname{EL}}
\newcommand{\U}{\mathcal U}
\newcommand{\Rad}{\operatorname{Rad}}
\newcommand{\tr}{\operatorname{tr}}
\newcommand{\opnorm}[1]{\left\lVert #1\right\rVert}
\newcommand{\hsnorm}[1]{\left\lVert #1\right\rVert_2}
\newcommand{\normal}[1]{\langle\!\langle #1\rangle\!\rangle}

\makeatletter
\patchcmd{\abstract}{\item[\hskip\labelsep\scshape\abstractname.]}{\item[]}{}{\PackageError{self-compression}{Could not remove abstract heading}{}}
\makeatother
\newcommand{\draftnotice}{%
  \begin{center}
    \small\itshape Draft. This manuscript is under revision.
  \end{center}}

\begin{document}
\frenchspacing
\title{Non-MF Groups}
\author{\paperauthor}
\subjclass[2020]{Primary 46L05; Secondary 20F65, 22D10, 16S50}
\keywords{MF groups, norm matrix coronas, property~(T), asymptotic representations, Leavitt algebra}

\begin{abstract}
\addtolength{\leftskip}{1.5pc}\addtolength{\rightskip}{1.5pc}
\leanverified{Manuscript/OneSidedMFRadical/CountableNonMF}{GroupApproximation.Manuscript.OneSidedMFRadical.manuscriptNotEveryCountableGroupIsMF}
Some groups are not MF\@.  The obstruction is a property-\textup{(T)}
subgroup conjugated into a proper subgroup of itself, together with an
element that commutes with the subgroup while its conjugate does not;
the matrices assigned to the resulting commutators converge to the
identity in normalized Hilbert--Schmidt norm, and on a normal subgroup
with property~\textup{(T)} inside the subgroup they generate, in
operator norm.
If a countable unital ring $R$ contains $s,t$ with $ts=1$ and
$R(1-st)R=R$, then every homomorphism from $\EL_n(R)$ to an MF group is
trivial for $n\ge2$.  In particular the unit group of the binary Leavitt
algebra $L_{\F_2}(1,2)$ is a finitely generated simple
property-\textup{(T)} group with no nontrivial MF image, so its reduced
$C^*$-algebra is separable and stably finite but not MF\@.  More
generally, for a countable purely infinite simple ring the largest MF
quotient of $\GL_n(R)$ is $K_1(R)$.  We also construct a sofic group
that is not MF and whose canonical trace on the maximal group
$C^*$-algebra is amenable but not quasidiagonal, and a torsion-free
group with no nontrivial homomorphism to an MF group.
\end{abstract}

\maketitle
\draftnotice

\section{Introduction}

A countable group $G$ is
\emph{MF}~\cite{CDE} if there are positive
integers $d_n$ and maps $V_n\colon G\to\U(d_n)$, with $V_n(1)=1$, such that
\[
 \opnorm{V_n(gh)-V_n(g)V_n(h)}\longrightarrow0
 \qquad(g,h\in G),
\]
and
\[
 \limsup_n\opnorm{V_n(g)-1}>0
 \qquad(g\in G\setminus\{1\}).
\]
For countable $G$, the second condition
may equivalently be required
along the full sequence with a constant independent of
$g$~\cite[Propositions~2 and~7]{Korchagin}.  The strong convergence
convention of~\cite{GaoEtAl,Schafhauser} also requires the models to
reproduce the operator norms of the left regular representation, so a group
that is not MF as defined here is not MF in that convention either.
Equivalently, $G$ embeds
in the unitary group of the norm matrix corona
\[
 \mathcal Q_{\mathbf d}
 =\prod_nM_{d_n}(\C)\big/\bigoplus_nM_{d_n}(\C).
\]
Here and below, products of $C^*$-algebras mean bounded products, and
$\bigoplus$ denotes the ideal of norm-null sequences.  We call a
homomorphism $G\to\U(\mathcal Q_{\mathbf d})$ a \emph{corona
homomorphism}.  A separable
$C^*$-algebra is MF if it embeds in a norm matrix corona~\cite{BK}.
Write
\[
 \operatorname{Rad}_{\mathrm{MF}}(G)
 =\bigcap\{\ker f:f\colon G\to M,\ M\text{ an MF group}\}
\]
for the \emph{MF radical} of $G$.  For countable $G$ the quotient
$G/\operatorname{Rad}_{\mathrm{MF}}(G)$ is again MF, because countably many
homomorphisms to MF groups separate the elements outside the radical and
the direct sums of their models along a common sequence of matrix sizes
remain asymptotically multiplicative; so it is the largest MF quotient
of $G$.

Some groups admit no such approximation.  We construct groups for which
every homomorphism to an MF group is trivial, as well as a sofic group
with a nontrivial element killed by every MF homomorphism.

Our main example is the group
\[
 H=L_{\F_2}(1,2)^\times,
\]
the unit group of the binary Leavitt algebra: $H$ is simple, and
every homomorphism from $H$ to an MF group is trivial.  The right module
$R_R$ over $R=L_{\F_2}(1,2)$ satisfies $R_R\cong R_R\oplus R_R$, and $H$
is its automorphism group; Khanh--Thanh show that
$H\cong\GL_n(R)=\EL_n(R)$ for every $n\ge2$~\cite[Proposition~4.2 and
Corollary~4.4]{KhanhThanh}.  OpenAI used this group to construct the
first nonsofic group~\cite[Chapter~3]{OAI}.

We use the group commutator convention $[g,h]=ghg^{-1}h^{-1}$.
The groups here fail to be MF because they contain a
property-\textup{(T)} subgroup that is conjugated into a proper subgroup
of itself, together with an element that commutes with the subgroup while
its conjugate does not.  Call $u\in G$ a \emph{compressor} of
$L\le G$ when $uLu^{-1}\le L$.
An element $c$ of the centralizer $C_G(L)$ commutes with $L$, so $ucu^{-1}$
commutes with $uLu^{-1}$ but need not commute with the rest of $L$.  The
normal subgroup of $G$ generated by the commutators of all such elements
$ucu^{-1}$ with elements of $L$ is
\begin{equation}\label{eq:intrinsic-defect}
 \mathfrak D_G(L)
 =\normal{
   [ucu^{-1},\ell]:
   uLu^{-1}\le L,\ c\in C_G(L),\ \ell\in L
 }_G.
\end{equation}
For every homomorphism $f\colon G\to H$,
\begin{equation}\label{eq:defect-functorial}
 f(\mathfrak D_G(L))\le\mathfrak D_H(f(L)),
\end{equation}
\leanverified{Manuscript/OneSidedMFRadical/PrintedDefectFunctorial}{GroupApproximation.Manuscript.OneSidedMFRadical.manuscriptPrintedDefectFunctorial}
since compressors, centralizing elements, and commutators of $L$ map to
compressors, centralizing elements, and commutators of $f(L)$, and
conjugates map to conjugates.

The criterion behind every example is Theorem~\ref{thm:compression-criterion}:
for a property-\textup{(T)} subgroup $L\le G$, every normal
property-\textup{(T)} subgroup of $G$ inside $\mathfrak D_G(L)$ dies in
every MF image.

Asymptotic multiplicativity in operator norm makes the conjugation maps
$\operatorname{Ad}(V_n(g))\colon x\mapsto V_n(g)xV_n(g)^*$ an asymptotically
multiplicative family of unitaries on $M_{d_n}(\C)$ with its normalized
Hilbert--Schmidt inner product.  Here
$\opnorm{\operatorname{Ad}(A)-\operatorname{Ad}(B)}\le2\opnorm{A-B}$, with
the norm on the left taken in the operator algebra of this Hilbert space.
The classes of these maps form a homomorphism of $G$ into the unitary
group of the norm matrix corona with coordinate sizes $d_n^2$.  Let $U$
be the class of $(\operatorname{Ad}(V_n(u)))$ and $P$ the image of the
Kazhdan projection of $L$ in that corona.  Then $U^*PU$ projects onto the vectors fixed by $u^{-1}Lu$, a
subgroup containing $L$ because $uLu^{-1}\le L$; so $U^*PU\le P$, and
stable finiteness of that corona forces $U^*PU=P$.  So conjugation by $u$ preserves the
Hilbert--Schmidt asymptotic commutant of $L$, and the matrices assigned to an element of
$\mathfrak D_G(L)$ tend to $1$ in Hilbert--Schmidt norm in every operator
norm asymptotic representation (Corollary~\ref{cor:defect-hs}).
Now let $K\trianglelefteq G$ have property~\textup{(T)} and lie in
$\mathfrak D_G(L)$; property~\textup{(T)} turns this into triviality in
operator norm.  A corona homomorphism nontrivial on $K$ compresses to a corner where
the Kazhdan projection of $K$ vanishes.  A limit of the normalized traces
of that corner is a tracial state taking the value $1$ on every element
of $K$, hence is the trivial character; but the trivial character takes
the value $1$ on the Kazhdan projection, which is $0$ there
(Theorem~\ref{thm:normal-kazhdan}).\footnote{A
Hilbert--Schmidt bound on the multiplicative defect does not give an
operator norm bound on the conjugation maps, so this argument does not
apply to sofic or hyperlinear approximations.}

The stable letter of an ascending HNN extension of a property-\textup{(T)}
group along a proper self-embedding is a compressor.  For a unital ring
$R$, write $e_{ij}(a)=1+aE_{ij}$ and let $\EL_n(R)$ be the subgroup of
$\GL_n(R)$ generated by the elements $e_{ij}(a)$, $i\ne j$.  A pair
$s,t\in R$ with $ts=1$ gives a compressor in the same way: an explicit
element $u\in\EL_4(R)$ conjugates $\EL_3(R)$ into itself by
$e_{ij}(a)\mapsto e_{ij}(sat)$, and if $1-st$ generates $R$ as a
two-sided ideal, then one commutator $[ucu^{-1},\ell]$ normally
generates all of $\EL_4(R)$.  This happens once, in
the universal unital associative ring
$\mathcal C=\mathbb Z\langle s_0,s_1,t_0,t_1\rangle/
(t_is_j-\delta_{ij}:i,j\in\{0,1\})$, and $\EL_4(\mathcal C)$ then maps
to every $\EL_n(R)$ below with normally generating image.

Theorem~\ref{thm:full-defect-ring} carries this to every countable
unital ring with such a pair $s,t$ and to every $n\ge2$: every
homomorphism from $\EL_n(R)$ to an MF group is trivial.

Let $L_{\F_2}(1,2)$ denote the binary
Leavitt algebra~\cite{Leavitt}, the unital $\F_2$-algebra generated by
$s_0,s_1,t_0,t_1$ subject to
\begin{equation}\label{eq:leavitt}
 t_i s_j=\delta_{ij},
 \qquad s_0t_0+s_1t_1=1.
\end{equation}

\begin{mainthm}[the unit group of the binary Leavitt algebra]\label{thm:headline}
Let $R=L_{\F_2}(1,2)$ and let $H=R^\times$ be its unit group.  Then
$H\cong\EL_4(R)$, and $H$ is finitely generated, nontrivial, simple, has
property~\textup{(T)}, and every homomorphism from $H$ to an MF group
is trivial.  So $H$ is not MF, and
$C^*_{\mathrm r}(H)$ is separable and stably finite but not MF, while
$C^*_{\max}(H)$ is not finite: it contains a proper isometry.
\leanverified{Manuscript/OneSidedMFRadical/UnitGroupHeadline}{GroupApproximation.Manuscript.OneSidedMFRadical.UnitGroupHeadline.manuscriptUnitGroupHeadline}
\end{mainthm}

For a countable purely infinite simple ring $R$ and $n\ge1$,
Theorem~\ref{thm:mf-quotient-units} identifies $\Rad_{\mathrm{MF}}(\GL_n(R))$
with the commutator subgroup, and the largest MF quotient of $\GL_n(R)$
with $K_1(R)$.  For
the unit group of $L_k(1,d)$ the quotient is $k^\times/(k^\times)^{d-1}$.

\begin{mainthm}[an amenable nonquasidiagonal trace]
\label{thm:amenable-trace}
There is a sofic group $W=W_0\rtimes\mathbb Z$, with $W_0$ locally residually
finite, that is not MF\@.  The canonical trace on $C^*_{\max}(W_0)$ is
quasidiagonal, and the canonical trace on $C^*_{\max}(W)$ is amenable
and not quasidiagonal.
\leanverified{Manuscript/OneSidedMFRadical/AmenableTraceTheorem}{GroupApproximation.AmenableTraceTheorem.manuscriptAmenableNonquasidiagonalTrace}
\end{mainthm}

Section~\ref{sec:amenable-nonqd} constructs $W$ from the affine group
$\mathbb Z^3\rtimes\mathrm{SL}_3(\mathbb Z)$ and Clifford lamps.\footnote{The
Lean proof that $W$ is not MF was registered at Palomar on August~24,
2026:
\url{https://palomar-registry.org/entry?id=PALOMAR-2026-08-24-000006&version=1}.}
The group $W$ is sofic because locally residually finite groups are sofic and
soficity passes to extensions with amenable
quotient~\cite[Theorem~1]{ElekSzabo}.  The trace answers the question of
Brown~\cite[discussion preceding Proposition~3.5.1]{Brown} and of
Schafhauser, Tikuisis and White~\cite[Problem~X(1)]{STW} whether every
amenable trace is quasidiagonal.
\leanverified{Manuscript/OneSidedMFRadical/AmenableTraceTheorem}{GroupApproximation.AmenableTraceTheorem.manuscriptWSoficFromLocallyRFExtension}
Tikuisis, White, and Winter proved
that faithful traces on separable nuclear $C^*$-algebras satisfying the
universal coefficient theorem are quasidiagonal~\cite{TWW}; here
$C^*_{\max}(W)$ is not nuclear, since $W$ is not amenable.  The canonical
trace of $C^*_{\max}(W_0)$ is quasidiagonal, so the crossed product by
$\mathbb Z$ destroys quasidiagonality of the canonical trace.  Since $W_0$ is a direct limit of
residually finite groups, it is MF~\cite[Corollary~10 and
Proposition~13]{Korchagin}.  So MF groups are not closed under semidirect
products with $\mathbb Z$, which answers a question of
Korchagin~\cite[remark following Proposition~12]{Korchagin}.
\leanverified{Manuscript/OneSidedMFRadical/AmenableTraceTheorem}{GroupApproximation.AmenableTraceTheorem.manuscriptShiftKernelIsOperatorMF}
\leanverified{Manuscript/OneSidedMFRadical/AmenableTraceTheorem}{GroupApproximation.AmenableTraceTheorem.manuscriptMFNotClosedUnderIntSemidirect}

Every example above has torsion, and in the lamp construction the
obstruction is itself a torsion element.  The last construction has
none.

\begin{mainthm}[a torsion-free example]
\label{thm:torsion-free}
There is a two-generated, finitely presented, torsion-free, acylindrically
hyperbolic group $Q$ with property~\textup{(T)} such that every
homomorphism from $Q$ to an MF group is trivial.  In particular, no
nontrivial quotient of $Q$ is MF\@.
\end{mainthm}

Section~\ref{sec:torsion-free} builds $Q$ from Fournier-Facio's
torsion-free property-\textup{(T)} group~\cite[\S2]{FFF} and Hull's small
cancellation theorem~\cite{Hull}.  Its reduced $C^*$-algebra is simple,
has a unique tracial state and stable rank one, and is not MF\@.

\subsection*{Related work}

Blackadar and Kirchberg introduced the notions of MF and NF algebras in a
1994 Oberwolfach abstract~\cite{BK94}.  They wrote that the classes of NF,
strong NF, and separable nuclear stably finite $C^*$-algebras might
coincide.  Their 1997 paper developed these notions and proved that a
separable $C^*$-algebra is NF if and only if it is nuclear and
MF~\cite{BK}.  It also asked whether every separable nuclear stably
finite $C^*$-algebra is quasidiagonal; for algebras that satisfy the
universal coefficient theorem and have a faithful trace, the theorem of
Tikuisis, White, and Winter recalled above answers this~\cite{TWW}.
The question without nuclearity, whether every stably finite $C^*$-algebra is MF, is the MF
problem, also posed by Blackadar and
Kirchberg~\cite[\S6]{GoldbringHart}.  A positive answer would make every
countable group MF, because the reduced group $C^*$-algebra is separable
and stably finite; Lubotzky--Oppenheim and, separately, Thom noted this when they
listed operator norm approximation of groups among the open
approximation problems~\cite{LubotzkyOppenheim,Thom}.  Carri\'on,
Dadarlat, and Eckhardt introduced MF groups~\cite{CDE}, and whether every
countable group is MF remained open~\cite{Korchagin,BDL}.  The MF problem
itself was answered negatively in 2020: the failure of the Connes
embedding problem~\cite{MIPRE} yields a $\mathrm{II}_1$ factor that
embeds in no norm ultraproduct of matrix algebras, hence a separable
stably finite $C^*$-algebra that is not MF~\cite[Proposition~6.1 and
Remark~6.2]{GoldbringHart}.  Theorem~\ref{thm:headline} gives a
counterexample among reduced group $C^*$-algebras and answers the group
question.

Theorem~\ref{thm:compression-criterion} uses the configuration of the
nonsofic construction~\cite[Proposition~2.3]{OAI}: a property-\textup{(T)}
subgroup $L$ conjugated into itself by an element $u$, and an element $c$
commuting with $L$ whose conjugate $ucu^{-1}$ does not.  There, soficity
and the rigidity theorems of Kun~\cite{Kun16} and Kun--Thom~\cite{KT19}
force the centralizing subgroup to be locally embeddable into finite
groups, and a copy of Thompson's group $V$ gives the contradiction.  Here
the rigidity comes from the Kazhdan projection of $L$ in a norm matrix
corona, which commutes with the compressor because the corona is stably
finite, and the contradiction is a commutator $[ucu^{-1},\ell]\ne1$
that every operator norm asymptotic representation sends to $1$ in
Hilbert--Schmidt norm.  The group $H$ is the group of~\cite{OAI}, and
Theorem~\ref{thm:torsion-free} reuses Fournier-Facio's
configuration~\cite{FFF}.  Bachner--Dogon--Lubotzky showed that for groups
of Deligne type, operator--Hilbert--Schmidt stability would imply that the
group is not MF~\cite[Proposition~1.5]{BDL}.  Here the compression
relation replaces the stability hypothesis.

\section{One-sided compression}
\label{sec:compression-radical}

\subsection{Corona homomorphisms}

\begin{lemma}\label{prop:mf-residual-calculus}
Let $G$ be countable and $K\le G$.  Every homomorphism from $G$ to an MF
group is trivial on $K$ if and only if every corona homomorphism from $G$
is trivial on $K$.  Let $A$ be a unital $C^*$-algebra and let
$\pi\colon G\to\U(A)$ be an injective group homomorphism.  If $A$
embeds in a norm matrix corona, then $G$ is MF\@.  In particular, this
applies to the canonical group homomorphisms into $\U(C^*_{\max}(G))$
and $\U(C^*_{\mathrm r}(G))$.
\leanverified{Analysis/MFAlgebra}{GroupApproximation.isCDEOperatorMF_of_faithful_corona_map}
\leanverified{Manuscript/OneSidedMFRadical/PrintedDefinitions}{GroupApproximation.Manuscript.OneSidedMFRadical.allMFTargetsKill_iff_allCoronasKill}
\end{lemma}

\begin{proof}
The image of a corona homomorphism is a countable MF group.  Conversely,
compose a homomorphism to an MF group with a corona embedding of its image;
the resulting corona homomorphism has the same kernel.  If $\iota$ is
an embedding as in the last assertion, then
$g\mapsto\iota(\pi(g))+1-\iota(1)$ is a corona homomorphism that is
injective on $G$; the complement term allows $\iota$ to be nonunital.
\end{proof}

\subsection{Kazhdan transport in normalized Hilbert--Schmidt norm}

For an actual representation $\rho\colon G\to\U(d)$ and $u\in G$ with
$uLu^{-1}\le L$, the commutant $\mathcal C=\rho(L)'\subseteq M_d(\C)$
satisfies $\rho(u)^*\mathcal C\rho(u)\subseteq\mathcal C$, and the two
spaces have the same finite dimension, so the inclusion is an equality;
if $c$ commutes with $L$, then $\rho(ucu^{-1})$ commutes with $\rho(L)$,
and $\rho$ is trivial on $\mathfrak D_G(L)$.  An asymptotic
representation has no exact commutant to count.
Property~\textup{(T)} supplies a Kazhdan projection in its place, and
stable finiteness of the corona replaces the dimension count.

For $a\in M_d(\C)$ write
\[
 \hsnorm{a}=\tr_d(a^*a)^{1/2},
 \qquad \tr_d=\frac1d\operatorname{Tr}.
\]
An \emph{operator norm asymptotic representation} of a group $G$
is a sequence of maps $V_n\colon G\to\U(d_n)$ such that $V_n(1)=1$ and
\[
 \opnorm{V_n(gh)-V_n(g)V_n(h)}\longrightarrow0
 \qquad(g,h\in G).
\]
The elements $g$ with $\hsnorm{V_n(g)-1}\to0$ form a normal subgroup of
$G$, because $\hsnorm{a}\le\opnorm{a}$, the Hilbert--Schmidt norm is
invariant under adjoints and unitary conjugation, and
$\opnorm{V_n(g^{-1})-V_n(g)^*}\to0$:
\begin{gather*}
 \hsnorm{V_n(gh)-1}\le\hsnorm{V_n(g)-1}+\hsnorm{V_n(h)-1}+o(1),\\
 \hsnorm{V_n(g^{-1})-1}=\hsnorm{V_n(g)-1}+o(1),\qquad
 \hsnorm{V_n(ghg^{-1})-1}=\hsnorm{V_n(h)-1}+o(1).
\end{gather*}
For $L\le G$, let
\[
 \mathcal C_2(V,L)
 =\left\{(x_n)\ \middle|\
   \begin{aligned}
    &\sup_n\hsnorm{x_n}<\infty,\\
    &\hsnorm{V_n(\ell)x_n-x_nV_n(\ell)}\longrightarrow0 &&(\ell\in L)
   \end{aligned}\right\}
\]
be the Hilbert--Schmidt asymptotic commutant of $V(L)$.  Coordinatewise
conjugation by $V_n(g)$ defines a bijection, denoted $\operatorname{Ad}(V(g))$,
of the Hilbert--Schmidt bounded matrix sequences.
A unital $C^*$-algebra is finite if each of its isometries is unitary,
and stably finite if all matrix algebras over it are finite.

\begin{lemma}\label{lem:stable-finite}
For every sequence of positive integers $(m_n)$, the norm matrix corona
\[
 \prod_nM_{m_n}(\C)\big/\bigoplus_nM_{m_n}(\C)
\]
is stably finite.  Consequently, unitarily equivalent projections in a norm
matrix corona are equal whenever one is dominated by the other.
\leanverified{Manuscript/OneSidedMFRadical/StableFiniteness}{GroupApproximation.Manuscript.OneSidedMFRadical.manuscriptNormMatrixCoronaStableFinite}
\end{lemma}

\begin{proof}
Write $A$ for the displayed corona.  If $v^*v=1$ in $A$ and $(x_n)$ is a
bounded lift of $v$, then $\opnorm{x_n^*x_n-1}\to0$, so for all large $n$
the matrix $u_n=x_n(x_n^*x_n)^{-1/2}$ is unitary and $\opnorm{u_n-x_n}\to0$;
after assigning arbitrary unitary values to the finitely many remaining
coordinates, $(u_n)$ is a unitary lift of $v$, and $vv^*=1$.  So $A$ is
finite, and so is $M_k(A)$, which is the corona with matrix sizes $km_n$.
Finally, if $p\le q$ are equivalent projections in a finite algebra
$B$ and $w$ is a partial isometry with $w^*w=q$ and $ww^*=p$, then
$w=qwq$, so $w+(1-q)$ is an isometry of $B$ and hence unitary; therefore
$p+(1-q)=(w+(1-q))(w+(1-q))^*=1$ and $p=q$.
\end{proof}

For a group $L$ with property~\textup{(T)}, the Kazhdan projection is the
central projection $e_L\in C^*_{\max}(L)$ whose image in every unitary
representation of $L$ is the orthogonal projection onto the
$L$-invariant vectors~\cite{AkemannWalter}.

\begin{lemma}[one-sided order for the Kazhdan projection]
\label{lem:kazhdan-projection-order}
Let $L$ have property~\textup{(T)}, let $B$ be a unital $C^*$-algebra, and
let $\pi\colon L\to\U(B)$ be a homomorphism.  Let $P\in B$ be the image
of $e_L$ under the $*$-homomorphism $C^*_{\max}(L)\to B$ induced by
$\pi$.  If a unitary $U\in B$ satisfies $U\pi(L)U^*\subseteq\pi(L)$, then
$U^*PU\le P$.
\leanverified{Manuscript/OneSidedMFRadical/KazhdanProjectionOrderLiteral}{GroupApproximation.Manuscript.OneSidedMFRadical.manuscriptLiteralMaximalCStarKazhdanProjectionOrder}
\end{lemma}

\begin{proof}
Fix a faithful nondegenerate representation of $B$ on a Hilbert space
$\mathcal H$, and let $\operatorname{Fix}\pi(L)$ be the space of vectors
fixed by $\pi(L)$.  Then $P$ is the orthogonal projection onto
$\operatorname{Fix}\pi(L)$.  If $\xi\in\operatorname{Fix}\pi(L)$ and
$\ell\in L$, then $\pi(\ell)U^*\xi=U^*(U\pi(\ell)U^*)\xi=U^*\xi$, because
$U\pi(\ell)U^*$ belongs to $\pi(L)$.  So
$U^*\operatorname{Fix}\pi(L)\subseteq\operatorname{Fix}\pi(L)$.  Since
$U^*PU$ is the orthogonal projection onto $U^*\operatorname{Fix}\pi(L)$,
it follows that $U^*PU\le P$.
\end{proof}

\begin{theorem}[one-sided Kazhdan transport]\label{thm:transport}
Let $L\le G$ have property~\textup{(T)}, and let $u\in G$ satisfy
$uLu^{-1}\le L$.  Let $(V_n)$ be an operator norm asymptotic
representation of $G$.  Then
\[
 \operatorname{Ad}(V(u))\bigl(\mathcal C_2(V,L)\bigr)
 =\mathcal C_2(V,L).
\]
\leanverified{Manuscript/OneSidedMFRadical/TransportPrintedRoute}{GroupApproximation.Manuscript.OneSidedMFRadical.TransportPrintedRoute.manuscriptPrintedTransportHS}
\end{theorem}

\begin{proof}
For unitaries $A,B$,
\[
 \opnorm{\operatorname{Ad}(A)-\operatorname{Ad}(B)}\le2\opnorm{A-B},
\]
where the norm on the left is the operator norm on the Hilbert space
$(M_{d_n}(\C),\hsnorm{\,\cdot\,})$, so the maps
$\operatorname{Ad}(V_n(g))$ are asymptotically multiplicative in
operator norm and
\[
 \widetilde\sigma(g)=[\operatorname{Ad}(V_n(g))]_n\in\U(\mathcal B),
 \qquad
 \mathcal B=\prod_nB(M_{d_n}(\C))\big/\bigoplus_nB(M_{d_n}(\C)),
\]
is a homomorphism, where $B(M_{d_n}(\C))$ denotes the operators on the
Hilbert--Schmidt Hilbert space.  The algebra $\mathcal B$ is a norm
matrix corona with coordinate sizes $d_n^2$ after a choice of matrix
units.  Let $P\in\mathcal B$ be the image of the Kazhdan projection of
$L$~\cite{AkemannWalter} under $C^*_{\max}(L)\to\mathcal B$, and lift
$P$ to orthogonal projections $P_n$ on
$(M_{d_n}(\C),\hsnorm{\,\cdot\,})$ by functional calculus.

A Hilbert--Schmidt bounded sequence $(x_n)$ lies in $\mathcal C_2(V,L)$
if and only if
\[
 \hsnorm{P_nx_n-x_n}\longrightarrow0 .
\]
In one direction, fix $\varepsilon>0$ and,
by density of the group algebra in $C^*_{\max}(L)$, a finite set
$F\subseteq L$ and scalars $(a_\ell)_{\ell\in F}$ with
\[
 \Bigl\|\sum_{\ell\in F}a_\ell u_\ell-e_L\Bigr\|<\varepsilon,
\]
where $u_\ell$ is the canonical unitary of $\ell$; the trivial
character, which sends $e_L$ to $1$, gives
$\bigl|\sum_{\ell\in F}a_\ell-1\bigr|<\varepsilon$.  The element
$\sum_{\ell\in F}a_\ell\widetilde\sigma(\ell)$ of $\mathcal B$ is within
$\varepsilon$ of $P$, so
\[
 \limsup_n\opnorm{P_n-\sum_{\ell\in F}a_\ell\operatorname{Ad}(V_n(\ell))}
 \le\varepsilon .
\]
If $(x_n)\in\mathcal C_2(V,L)$ and $\sup_n\hsnorm{x_n}\le c$, then
$\operatorname{Ad}(V_n(\ell))x_n=x_n+o(1)$ in $\hsnorm{\,\cdot\,}$ for
each $\ell\in F$, so
\[
 \limsup_n\hsnorm{P_nx_n-x_n}
 \le\varepsilon c+\Bigl|\sum_{\ell\in F}a_\ell-1\Bigr|c
 \le2\varepsilon c,
\]
and $\varepsilon$ was arbitrary.  In the other direction,
$u_\ell e_L=e_L$ in $C^*_{\max}(L)$, so $\widetilde\sigma(\ell)P=P$ and
$\opnorm{\operatorname{Ad}(V_n(\ell))P_n-P_n}\to0$; if
$\hsnorm{P_nx_n-x_n}\to0$, then
\[
\begin{aligned}
 \hsnorm{\operatorname{Ad}(V_n(\ell))x_n-x_n}
 &\le2\hsnorm{x_n-P_nx_n}\\
 &\quad+\opnorm{\operatorname{Ad}(V_n(\ell))P_n-P_n}\,\hsnorm{x_n}
 \longrightarrow0 .
\end{aligned}
\]

Put $U=\widetilde\sigma(u)$.  Since $uLu^{-1}\le L$,
$U\widetilde\sigma(L)U^*\subseteq\widetilde\sigma(L)$, so $U^*PU\le P$
by Lemma~\ref{lem:kazhdan-projection-order}, and $U^*PU=P$ by
Lemma~\ref{lem:stable-finite}, since the two projections are unitarily
equivalent.  So $UP=PU$, and correspondingly
$\opnorm{\operatorname{Ad}(V_n(u))P_n-P_n\operatorname{Ad}(V_n(u))}\to0$.  For
$(x_n)\in\mathcal C_2(V,L)$,
\[
\begin{aligned}
 \hsnorm{P_n\operatorname{Ad}(V_n(u))^{\pm1}x_n
 -\operatorname{Ad}(V_n(u))^{\pm1}x_n}
 &\le\opnorm{[P_n,\operatorname{Ad}(V_n(u))^{\pm1}]}\,\hsnorm{x_n}\\
 &\quad+\hsnorm{P_nx_n-x_n}\longrightarrow0,
\end{aligned}
\]
so $\operatorname{Ad}(V(u))^{\pm1}x\in\mathcal C_2(V,L)$ by the
displayed equivalence.
\end{proof}

\begin{corollary}\label{cor:defect-hs}
Let $L\le G$ have property~\textup{(T)} and let $(V_n)$ be an operator
norm asymptotic representation of $G$.  Then
\[
 \hsnorm{V_n(d)-1}\longrightarrow0
 \qquad(d\in\mathfrak D_G(L)).
\]
\leanverified{Manuscript/OneSidedMFRadical/DefectHS}{GroupApproximation.Manuscript.OneSidedMFRadical.manuscriptCompressionDefectHSInvisible}
\leanverified{Manuscript/OneSidedMFRadical/PrintedDefectShadow}{GroupApproximation.Manuscript.OneSidedMFRadical.manuscriptPrintedDefectShadowInclusion}
\end{corollary}

\begin{proof}
Let $uLu^{-1}\le L$, $c\in C_G(L)$, and $\ell\in L$, and
write $\mathcal C_2=\mathcal C_2(V,L)$.  Since $c$ commutes with $L$,
$(V_n(c))\in\mathcal C_2$, so $(V_n(u)V_n(c)V_n(u)^*)\in\mathcal C_2$ by
Theorem~\ref{thm:transport}, and then $(V_n(ucu^{-1}))\in\mathcal C_2$
because $\opnorm{V_n(u)V_n(c)V_n(u)^*-V_n(ucu^{-1})}\to0$.  So
$\hsnorm{V_n(\ell)V_n(ucu^{-1})-V_n(ucu^{-1})V_n(\ell)}\to0$, and by
asymptotic multiplicativity $\hsnorm{V_n([ucu^{-1},\ell])-1}\to0$.  The
elements with this property form a normal subgroup of $G$, so it
contains $\mathfrak D_G(L)$.
\end{proof}

\subsection{From Hilbert--Schmidt to operator norm}

Smallness in normalized Hilbert--Schmidt norm does not bound the
operator norm: for
\[
 D_d=\operatorname{diag}(-1,1,\ldots,1)\in\U(d),
 \qquad
 \hsnorm{D_d-1}=\frac{2}{\sqrt d}\longrightarrow0,
 \qquad
 \opnorm{D_d-1}=2 .
\]
A corner occupying a vanishing fraction of the matrix can therefore
carry the entire obstruction, which is why the corner below is
renormalized by its own rank.

Restricting a corona homomorphism $\rho$ to a corner requires a
correction: a projection $q$ commuting with $\rho(G)$ has lifts $q_n$
that commute only asymptotically with unitary lifts $U_n(g)$ of
$\rho(g)$, so the compressions $q_nU_n(g)q_n$ are only approximately
unitary in the matrix corners.

\begin{lemma}[central corona corners]\label{lem:central-corona-corner}
Let \(\rho\colon G\to\U(\mathcal Q_{\mathbf d})\) be a homomorphism from a
countable group, and let \(q\in\mathcal Q_{\mathbf d}\) be a nonzero
projection commuting with \(\rho(G)\).  After passing to an infinite
subsequence of coordinates, there are ranks \(r_n\ge1\), an
identification of the corner \(q\mathcal Q_{\mathbf d}q\) with the norm
matrix corona \(\mathcal Q_{\mathbf r}\), and an operator norm asymptotic
representation \(W_n\colon G\to\U(r_n)\) whose corona homomorphism is
\(g\mapsto q\rho(g)\) under this identification.
\leanverified{Manuscript/OneSidedMFRadical/CentralCoronaCorner}{GroupApproximation.Manuscript.OneSidedMFRadical.manuscriptCentralCoronaCorner}
\leanverified{Manuscript/OneSidedMFRadical/CentralCoronaCornerPrintedRoute}{GroupApproximation.Manuscript.OneSidedMFRadical.CentralCoronaCornerPrintedRoute.manuscriptPrintedCentralCoronaCorner}
\end{lemma}

\begin{proof}
Lift \(q\) to projections \(q_n\in M_{d_n}(\C)\) by functional calculus.
Since \(q\ne0\), infinitely many \(q_n\) are nonzero; retain those
coordinates and put \(r_n=\operatorname{rank}(q_n)\).  Identifying each corner \(q_nM_{d_n}(\C)q_n\) with \(M_{r_n}(\C)\)
identifies the corner
\(q\mathcal Q_{\mathbf d}q\) with the corona \(\mathcal Q_{\mathbf r}\):
a bounded sequence \((z_n)\) representing \(z\in\mathcal Q_{\mathbf d}\)
gives the representative \((q_nz_nq_n)\) of \(qzq\).  Since \(q\)
commutes with \(\rho(G)\), the map \(g\mapsto q\rho(g)\) is a
homomorphism from \(G\) to the unitary group of the corner, with unit
\(q\).  Each value is a unitary of \(\mathcal Q_{\mathbf r}\), so, as in
the proof of Lemma~\ref{lem:stable-finite}, it lifts to unitaries
\(W_n(g)\in\U(r_n)\), with \(W_n(1)=I_{r_n}\).  Multiplicativity in the
corona makes \((W_n)\) an operator norm asymptotic representation whose
corona homomorphism is \(g\mapsto q\rho(g)\).
\end{proof}

\begin{theorem}[normal Kazhdan subgroups]\label{thm:normal-kazhdan}
Let \(G\) be countable and let \(K\trianglelefteq G\) have
property~\textup{(T)}.  If every operator norm asymptotic representation
\((V_n)\) of \(G\) satisfies \(\hsnorm{V_n(k)-1}\to0\) for all \(k\in K\),
then every corona homomorphism from \(G\) is trivial on \(K\).
\leanverified{Manuscript/OneSidedMFRadical/NormalKazhdan}{GroupApproximation.Manuscript.OneSidedMFRadical.manuscriptNormalKazhdanRadical}
\leanverified{Manuscript/OneSidedMFRadical/NormalKazhdanPrintedRoute}{GroupApproximation.Manuscript.OneSidedMFRadical.NormalKazhdanPrintedRoute.manuscriptNormalKazhdanRadical_printedRoute}
\end{theorem}

\begin{proof}
Suppose that a corona homomorphism
\(\Theta\colon G\to\U(\mathcal Q_{\mathbf d})\) is nontrivial on \(K\).
Let \(e_K\in C^*_{\max}(K)\) be the Kazhdan projection of \(K\), and let
\(p\in\mathcal Q_{\mathbf d}\) be its image under the homomorphism
\(C^*_{\max}(K)\to\mathcal Q_{\mathbf d}\) induced by \(\Theta|_K\).  In
any faithful representation of the corona, \(p\) is the projection onto
the \(K\)-fixed vectors.  Since \(gKg^{-1}=K\), the range of
\(\Theta(g)p\Theta(g)^*\) is
\(\Theta(g)\operatorname{Fix}\Theta(K)
=\operatorname{Fix}\Theta(gKg^{-1})=\operatorname{Fix}\Theta(K)\), so
\[
 \Theta(g)p\Theta(g)^*=p\qquad(g\in G),
\]
and \(q=1-p\) is nonzero because \(\Theta\) is nontrivial on \(K\).  By
Lemma~\ref{lem:central-corona-corner} there is an operator norm
asymptotic representation \(W_n\colon G\to\U(r_n)\) whose corona
homomorphism \(\widehat\Theta\) is the coordinate restriction of
\(g\mapsto q\Theta(g)\).  Let
\(\pi\colon C^*_{\max}(K)\to\mathcal Q_{\mathbf r}\) be the homomorphism
induced by \(\widehat\Theta|_K\).  Since \(q\) commutes with \(\Theta(K)\),
the map \(a\mapsto q\Theta(a)\) is a homomorphism on \(C^*_{\max}(K)\)
agreeing with \(\pi\) on \(K\), so \(\pi(e_K)\) is the coordinate
restriction of \(qp=0\).

Fix a free ultrafilter \(\omega\) and let \(\tau\) be the limit along
\(\omega\) of the normalized traces of the coordinates \(M_{r_n}(\C)\).
A norm-null sequence has vanishing traces, since
\(|\tr_{r_n}(x)|\le\opnorm x\), so \(\tau\) is a well defined tracial
state on \(\mathcal Q_{\mathbf r}\).  The hypothesis applies to
\((W_n)\) itself, so
\[
 |\tr_{r_n}(W_n(k))-1|\le\hsnorm{W_n(k)-I_{r_n}}\longrightarrow0
 \qquad(k\in K)
\]
in the normalized Hilbert--Schmidt norm of \(M_{r_n}(\C)\), and
therefore \(\tau(\pi(u_k))=1\) for every \(k\in K\), where \(u_k\) is
the canonical unitary of \(k\).  The trivial character \(\chi\) of
\(K\) is the state of \(C^*_{\max}(K)\) with \(\chi(u_k)=1\) for every
\(k\), and it has \(\chi(e_K)=1\).  The two states \(\tau\circ\pi\) and
\(\chi\) agree on every \(u_k\), hence by linearity on the dense
subalgebra \(\C[K]\), hence on \(C^*_{\max}(K)\).  So
\[
 1=\chi(e_K)=\tau(\pi(e_K))=\tau(0)=0,
\]
a contradiction.
\end{proof}

\begin{theorem}[one-sided compression criterion]
\label{thm:compression-criterion}
Let $G$ be countable and let $L\le G$ have property~\textup{(T)}.  If
$K\trianglelefteq G$ has property~\textup{(T)} and
$K\le\mathfrak D_G(L)$, then every homomorphism from $G$ to an MF group is
trivial on $K$.  In particular, if such a $K$ is nontrivial, then $G$ is
not MF, and
if $G$ has property~\textup{(T)} and $\mathfrak D_G(L)=G$, then every
homomorphism from $G$ to an MF group is trivial.
\leanverified{Manuscript/OneSidedMFRadical/PrintedCriterion}{GroupApproximation.Manuscript.OneSidedMFRadical.manuscriptOneSidedCompressionCriterion}
\leanverified{Manuscript/OneSidedMFRadical/PrintedCriterion}{GroupApproximation.Manuscript.OneSidedMFRadical.manuscriptFullRadicalKillsMFTargets}
\end{theorem}

\begin{proof}
By Corollary~\ref{cor:defect-hs}, every operator norm asymptotic
representation of \(G\) satisfies \(\hsnorm{V_n(k)-1}\to0\) for
\(k\in K\le\mathfrak D_G(L)\).  By Theorem~\ref{thm:normal-kazhdan}, every
corona homomorphism is then trivial on \(K\), and by
Lemma~\ref{prop:mf-residual-calculus} so is every homomorphism to an MF
group.  If \(K\ne1\), then \(G\) does not embed in an MF group, so \(G\)
is not MF\@.  The last assertion is the case \(K=G=\mathfrak D_G(L)\).
\end{proof}

\subsection{The maximal group \texorpdfstring{$C^*$}{C*}-algebra}

\begin{proposition}[proper one-sided compression]\label{prop:max-infinite}
Let $G$ be a countable group containing a property-\textup{(T)} subgroup
$\Gamma$ and an element $t$ with $t\Gamma t^{-1}\subsetneq\Gamma$.  Then
$C^*_{\max}(G)$ contains a proper isometry.  In particular,
$C^*_{\max}(G)$ is not stably finite, admits no faithful tracial state,
and is neither residually finite-dimensional nor MF\@.
\leanverified{Analysis/StrictCompressionFromPrinted}{GroupApproximation.MaximalCStarPrintedHypotheses.manuscriptMaximalCStarRemarkFromPrintedHypotheses}
\leanverified{Analysis/QuasiRegularWitness}{GroupApproximation.QuasiRegularWitness.baseVector_apply_base}
\leanverified{Analysis/ProperIsometryStrictOrder}{GroupApproximation.manuscriptProperIsometryStrictOrder}
\end{proposition}

\begin{proof}
Let $P\in C^*_{\max}(G)$ be the image of the Kazhdan projection of
$\Gamma$~\cite{AkemannWalter} under the canonical map
$C^*_{\max}(\Gamma)\to C^*_{\max}(G)$, and let $u$ be the canonical
unitary of $t$.  Conjugation by $u$ implements the isomorphism
$\Gamma\to t\Gamma t^{-1}$ on canonical unitaries, so by
Lemma~\ref{lem:kazhdan-projection-order} applied with $U=u$,
$P\le uPu^{*}$, where $uPu^{*}$ is the image of the Kazhdan projection
of $t\Gamma t^{-1}$.  The inequality is strict: in the quasi-regular
representation of $G$ on $\ell^2(G/t\Gamma t^{-1})$, the point mass at
the base coset is fixed by $t\Gamma t^{-1}$ and moved by every element of
$\Gamma\setminus t\Gamma t^{-1}$, so the two fixed-space projections
differ.  Put $q=uPu^*$, so that $P<q$, $Pq=qP=P$, $uP=qu$, and
$Pu^*(1-q)=Pu^*-Pu^*uPu^*=0$, and put $s=Pu^*+(1-q)$.  Then
\[
 s^*s=uPu^*+(1-q)=1,\qquad ss^*=P+(1-q)=1-(q-P)\ne1,
\]
so $s$ is a proper isometry, and $\operatorname{diag}(s,1,\ldots,1)$
is one in every matrix algebra over $C^*_{\max}(G)$, which is
therefore not stably finite.  A tracial state $\tau$ has
$\tau(q-P)=\tau(s^*s)-\tau(ss^*)=0$ with $q-P$ a nonzero positive
element, so no tracial state is faithful.  MF algebras are stably
finite, and a separable residually finite-dimensional $C^*$-algebra is
MF~\cite{BK}.
\end{proof}

A group can satisfy the hypothesis of
Proposition~\ref{prop:max-infinite} and be MF\@.  The ascending HNN
extension $V$ of $\mathbb Z^3\rtimes\mathrm{SL}_3(\mathbb Z)$ along
$(v,A)\mapsto(2v,A)$, used in Section~\ref{sec:amenable-nonqd}, has the
faithful matrix realization
\[
 V\cong
 \left\{
 \begin{pmatrix}
 2^kA&v\\0&1
 \end{pmatrix}
 :A\in\mathrm{SL}_3(\mathbb Z),\
 v\in\mathbb Z[1/2]^3,\ k\in\mathbb Z
 \right\},
\]
and reduction modulo odd integers separates its elements, so $V$ is
residually finite and MF~\cite[Corollary~10]{Korchagin}.  So the lamps
of Section~\ref{sec:amenable-nonqd} are necessary for the non-MF
conclusion there.

\section{One-sided inverses and elementary groups}
\label{sec:one-sided-inverses}

The proof of Theorem~\ref{thm:full-defect-ring} runs in rank four
over the ring $\mathcal C$, with three coordinates for the subgroup
$\EL_3(\mathcal C)$, which has property~\textup{(T)} because
$\mathcal C$ is finitely generated, and one for its centralizer; two
lemmas then carry the conclusion from $\EL_4(\mathcal C)$ to every
$\EL_n(R)$.  The
compressor needs only $ts=1$; the ideal condition on $1-st$ enters in
the last step of the rank-four argument and in the first lemma.

\begin{lemma}[two copies]\label{lem:two-copies}
Let $R$ be a unital ring with $s,t\in R$ such that $ts=1$ and
$R(1-st)R=R$.  Then there are $v_0,v_1,w_0,w_1\in R$ with
$w_iv_j=\delta_{ij}$.  Conversely, such elements satisfy $w_0v_0=1$ and
$w_1(1-v_0w_0)v_1=1$.
\leanverified{Manuscript/OneSidedMFRadical/RankDescentPrintedLemmas}{GroupApproximation.Manuscript.OneSidedMFRadical.RankDescentPrinted.manuscriptTwoCopiesLemma}
\end{lemma}

\begin{proof}
Put $e=1-st$, so that $es=te=0$, and choose $a_j,b_j\in R$, $0\le j<m$,
with $\sum_ja_jeb_j=1$.  Put
\[
 v_0=s^m,\qquad w_0=t^m,\qquad v_1=\sum_js^jeb_j,\qquad w_1=\sum_ja_jet^j.
\]
For all $i,j\ge0$, $et^is^je=\delta_{ij}e$: for $j>i$ the middle factor is
$s^{j-i}$, which $e$ kills on the left, and for $i>j$ it is $t^{i-j}$,
which $e$ kills on the right.  So $w_0v_0=t^ms^m=1$,
$w_1v_1=\sum_ja_jeb_j=1$, $w_0v_1=\sum_jt^{m-j}eb_j=0$, and
$w_1v_0=\sum_ja_jes^{m-j}=0$.  Conversely,
$w_1(1-v_0w_0)v_1=w_1v_1-(w_1v_0)(w_0v_1)=1$.
\end{proof}

\begin{lemma}[normal generation in rank two]\label{lem:rank-two}
Let $R$ be a unital ring and let $v,w,a,b\in R$ satisfy $wv=1$, $ba=1$,
and $bv=0$.  Then
\[
 D=[e_{12}(v),e_{21}(b)]=\operatorname{diag}(1+vb,1)
\]
normally generates $\EL_2(R)$.
\leanverified{Manuscript/OneSidedMFRadical/RankDescentPrintedLemmas}{GroupApproximation.Manuscript.OneSidedMFRadical.RankDescentPrinted.manuscriptRankTwoNormalGeneration}
\end{lemma}

\begin{proof}
The matrix $e_{12}(v)e_{21}(b)e_{12}(-v)$ has rows $(1+vb,\,-vbv)$ and
$(b,\,1-bv)$, so for $bv=0$ it is
$\begin{psmallmatrix}1+vb&0\\b&1\end{psmallmatrix}$, and multiplying by
$e_{21}(-b)$ gives $D=\operatorname{diag}(1+vb,1)$.  Let $N$ be the
normal closure of $D$ in $\EL_2(R)$.  For $r\in R$,
$[D,e_{12}(ar)]=e_{12}((1+vb)ar-ar)=e_{12}(vr)$, since $ba=1$; so
$e_{12}(vR)\le N$.  Put $f=1-vw$, an idempotent with $fv=0=wf$, and
\[
 z=vf+fw+1-f-vfw .
\]
Put $x=vf$, $y=fw$, $q=vfw$, and $c=1-f-q$, so that $z=x+y+c$.  From
$wv=1$ and $fv=wf=0$: $x^2=y^2=0$, $xy=q$, $yx=f$, the idempotents $f$
and $q$ are orthogonal, and $c$ annihilates $x$ and $y$ on both sides;
so $z^2=q+f+c=1$ and $zf=xf+yf+cf=vf$.  The factorization
\[
 \operatorname{diag}(z,z^{-1})
 =e_{12}(z)e_{21}(-z^{-1})e_{12}(z)\,e_{12}(-1)e_{21}(1)e_{12}(-1)
\]
puts $h=\operatorname{diag}(z,z)$ in $\EL_2(R)$, and
$he_{12}(fr)h^{-1}=e_{12}(zfrz)=e_{12}(v\,frz)\in N$.  So
$e_{12}(r)=e_{12}(vwr)\,e_{12}(fr)\in N$ for every $r\in R$, and
conjugating by $e_{12}(1)e_{21}(-1)e_{12}(1)$ gives $e_{21}(R)\le N$.
\end{proof}

\begin{lemma}[a rank-four compression cell]
\label{lem:ring-compression-cell}
Let $R$ be a unital ring, let $s,t\in R$ satisfy $ts=1$, and put
$e=1-st$.  In $G=\EL_4(R)$ let $L=\EL_3(R)$ occupy coordinates
$1,2,3$.  There are $u,c\in G$ with $uLu^{-1}\le L$ and $c\in C_G(L)$
such that
\[
 ucu^{-1}=e_{12}(e),
 \qquad
 [ucu^{-1},e_{23}(1)]=e_{13}(e).
\]
Every off-diagonal entry of a matrix in $uLu^{-1}$ has the form $sat$
with $a\in R$.
\end{lemma}

\begin{proof}
Here $e^2=e$ and $es=te=0$.  For $i=1,2,3$, set
\[
 u_i=e_{4i}(t-1)e_{i4}(1)e_{4i}(s-1)e_{i4}(-t).
\]
Its block on coordinates $(i,4)$ is
$\begin{psmallmatrix}s&e\\0&t\end{psmallmatrix}$, so
$u=u_3u_2u_1\in\EL_4(R)$ is the matrix
\[
 u=
 \begin{pmatrix}
 s&0&0&e\\
 0&s&0&et\\
 0&0&s&et^2\\
 0&0&0&t^3
 \end{pmatrix},
\]
invertible as a product of elementary matrices.  For $1\le i\ne j\le3$
and $a\in R$,
\begin{equation}\label{eq:intertwine}
 u\,e_{ij}(a)=e_{ij}(sat)\,u:
\end{equation}
both sides differ from $u$ by one entry, $sa$ in position $(i,j)$, since
on the right the increment $sat\cdot(s,\ldots,et^{j-1})$ from row $j$ of
$u$ contributes $sa\,ts=sa$ in position $(i,j)$ and
$sa\,te\,t^{j-1}=0$ in position $(i,4)$.  So $uLu^{-1}\le L$, and
\eqref{eq:intertwine} gives the last assertion.

The element
\[
 c=[e_{41}(e),e_{14}(t)]
   =\operatorname{diag}(1,1,1,1+et)
\]
is computed in the block on coordinates $(1,4)$ using $te=0$, and a
diagonal matrix of this shape commutes with every
$\operatorname{diag}(A,1)$, so $c\in C_G(L)$.  Both $uc$ and
$e_{12}(e)u$ equal $u+et\,E_{14}$: the last column of $uc$ is
$(e,et,et^2,t^3)^{\mathsf T}(1+et)$ and $e^2=e$, $te=0$, while
$e_{12}(e)u$ adds $e$ times the second row $(0,s,0,et)$ of $u$ to
its first row, and $es=0$.  So
$ucu^{-1}=e_{12}(e)$, and the Steinberg relation
$[e_{12}(e),e_{23}(1)]=e_{13}(e)$ gives the second identity.
\end{proof}

% The full-defect theorem and its rank-two descent have the selected
% Lean counterparts listed at Theorem~\ref{thm:full-defect-ring}.
% Sentence-level coverage is recorded separately in the sentence census.
\begin{theorem}[full complementary idempotents]\label{thm:full-defect-ring}
Let $R$ be a countable unital associative ring.  Suppose that $s,t\in R$
satisfy
\[
 ts=1,\qquad R(1-st)R=R;
\]
that is, $1-st$ is a \emph{full} idempotent, so that
$1=\sum_{j}a_j(1-st)b_j$ for finitely many $a_j,b_j\in R$.
For every $n\ge2$, every homomorphism from $\EL_n(R)$ to an MF group is
trivial.  Every homomorphism from $\EL_4(\mathcal C)$ to an MF group is
trivial; that group is finitely generated and has
property~\textup{(T)}; and there is a homomorphism from it to
$\EL_n(R)$ whose image normally generates $\EL_n(R)$.
\leanverified{Manuscript/OneSidedMFRadical/FullDefectRingEJZUnconditional}{GroupApproximation.Manuscript.OneSidedMFRadical.FullDefectRingEJZUnconditional.manuscriptFullComplementaryIdempotentsRankTwoAllCharacteristics}
\leanverified{Manuscript/OneSidedMFRadical/UniversalGroupSigma}{GroupApproximation.Manuscript.OneSidedMFRadical.UniversalGroupSigma.manuscriptFullComplementaryIdempotentsUniversal}
\end{theorem}

\begin{proof}
\emph{The universal group.}  In $\mathcal C$ put $s=s_0$, $t=t_0$, and
$e=1-st$.  Then $ts=1$, $e^2=e$, $es=te=0$, and $t_1es_1=1$.  Put
$G=\EL_4(\mathcal C)$ and $L=\EL_3(\mathcal C)$ on coordinates
$1,2,3$.  The ring $\mathcal C$ is nonzero, since it maps onto
$L_{\F_2}(1,2)$, and finitely generated, so both groups
have property~\textup{(T)} by~\cite[Theorem~1.1]{EJZ}, and
$\EL_4(\mathcal C)$ is finitely
generated~\cite[Theorem~1.3.1]{BHV}.

By Lemma~\ref{lem:ring-compression-cell} there are $u,c\in G$ with
$uLu^{-1}\le L$, $c\in C_G(L)$, and
\[
 d=[ucu^{-1},e_{23}(1)]=e_{13}(e)\in\mathfrak D_G(L).
\]

Let $N$ be the normal closure of $d$ in $G$.  For arbitrary $a,b\in\mathcal C$,
the Steinberg relations give
\[
 [e_{41}(a),e_{13}(e)]=e_{43}(ae),\qquad
 [e_{43}(ae),e_{32}(b)]=e_{42}(aeb).
\]
So $e_{42}(1)=e_{42}(t_1es_1)\in N$.
Elementary signed permutation matrices conjugate this element to every
off-diagonal position, up to a sign, so $e_{ij}(1)\in N$ whenever $i\ne j$.
For distinct $i,j,k$ and every $r\in\mathcal C$,
\[
 [e_{ij}(1),e_{jk}(r)]=e_{ik}(r)\in N.
\]
So $N=G$, and $\mathfrak D_G(L)=G$.  By
Theorem~\ref{thm:compression-criterion} with $K=G$, every homomorphism
from $G$ to an MF group is trivial.

\emph{Rank two.}  Now let $R$ satisfy the hypothesis.  By
Lemma~\ref{lem:two-copies} choose $v_0,v_1,w_0,w_1\in R$ with
$w_iv_j=\delta_{ij}$.  Then $\varphi\colon s_i\mapsto v_i$, $t_i\mapsto w_i$ is a unital ring
homomorphism $\mathcal C\to R$, and it induces a homomorphism
$\EL_4(\mathcal C)\to\EL_4(R)$.  Put
\begin{gather*}
 S_1=v_0v_0,\quad S_2=v_0v_1,\quad S_3=v_1v_0,\quad S_4=v_1v_1,\\
 T_1=w_0w_0,\quad T_2=w_1w_0,\quad T_3=w_0w_1,\quad T_4=w_1w_1,
\end{gather*}
so that $T_iS_j=\delta_{ij}$, and put $p=\sum_iS_iT_i$.  For a
$4\times4$ matrix $A$ over $R$ put
\[
 \jmath(A)=1-p+\sum_{i,j}S_iA_{ij}T_j\in R .
\]
Since $(1-p)S_i=0$ and $T_j(1-p)=0$, we have
$\jmath(AB)=\jmath(A)\jmath(B)$ and $\jmath(I)=1$, so
$\Psi(A)=\operatorname{diag}(\jmath(A),1)$ defines a homomorphism
$\GL_4(R)\to\GL_2(R)$, and for $i\ne j$ and $r\in R$,
\[
 \Psi(e_{ij}(r))=\operatorname{diag}(1+S_irT_j,1)=[e_{12}(S_ir),e_{21}(T_j)]
\]
by the computation in Lemma~\ref{lem:rank-two}, since $T_jS_i=0$.  So
$\Psi(\EL_4(R))\le\EL_2(R)$.  Let $\rho$ be a homomorphism from
$\EL_2(R)$ to an MF group.  Composed with $\EL_4(\mathcal C)\to\EL_4(R)\to\EL_2(R)$
it is trivial, so $\rho$ kills the image
$\Psi(e_{12}(1))=\operatorname{diag}(1+S_1T_2,1)$ of
$e_{12}(1)\in\EL_4(\mathcal C)$,
which normally generates $\EL_2(R)$ by Lemma~\ref{lem:rank-two} with
$v=S_1$, $w=T_1$, $a=S_2$, $b=T_2$.  So $\rho$ is trivial, and the image
of $\EL_4(\mathcal C)$ in $\EL_2(R)$ normally generates $\EL_2(R)$.

\emph{All ranks.}  For $n\ge3$, signed permutation matrices conjugate
$e_{12}(r)$ to $e_{ij}(\pm r)$ for all $i\ne j$, so the copy of
$\EL_2(R)$ on coordinates $1,2$ normally generates $\EL_n(R)$.  A
homomorphism from $\EL_n(R)$ to an MF group is trivial on that copy,
hence trivial, and the image of $\EL_4(\mathcal C)$ in that copy
normally generates
$\EL_n(R)$.
\end{proof}

\begin{corollary}\label{cor:simple-infinite-ring}
If $R$ is a countable simple unital ring that is not directly finite,
then every homomorphism from $\EL_n(R)$ to an MF group is trivial for
every $n\ge2$.
The same conclusion holds for $R=L_k(1,m)$, for every countable field
$k$ and every $m\ge2$.
\leanverified{Manuscript/OneSidedMFRadical/FullDefectRingEJZUnconditional}{GroupApproximation.Manuscript.OneSidedMFRadical.FullDefectRingEJZUnconditional.manuscriptSimpleInfiniteRingRankTwoAllCharacteristics}
\leanverified{Manuscript/OneSidedMFRadical/FullDefectRingEJZUnconditional}{GroupApproximation.Manuscript.OneSidedMFRadical.FullDefectRingEJZUnconditional.manuscriptLeavittAlgebraFullDefectRankTwoAllCharacteristics}
\end{corollary}

\begin{proof}
In the first case choose $ts=1\ne st$; then $1-st$ is a nonzero
idempotent, and since $R$ is simple it generates $R$ as a two-sided
ideal.  In the second case, with generators $s_1,\dots,s_m,t_1,\dots,t_m$
subject to $t_is_j=\delta_{ij}$ and $\sum_is_it_i=1$, take $s=s_1$ and
$t=t_1$; then $1-s_1t_1=\sum_{i\ge2}s_it_i$ and
\[
 t_2(1-s_1t_1)s_2=1.
\]
Theorem~\ref{thm:full-defect-ring} applies in both cases.
\end{proof}

\begin{corollary}\label{cor:one-sided-ring-maximal}
Let $R$ be a countable unital ring.  If $R$ is not directly finite,
then $C^*_{\max}(\EL_n(R))$ contains a proper isometry for every
$n\ge4$.  If $R\ne0$ satisfies the hypothesis of
Theorem~\ref{thm:full-defect-ring}, then for every $n\ge2$ the algebra
$C^*_{\mathrm r}(\EL_n(R))$ is separable, stably finite, and not MF,
and the unit group $R^\times$ is not MF\@.
\leanverified{Manuscript/OneSidedMFRadical/FullDefectRingEJZUnconditional}{GroupApproximation.Manuscript.OneSidedMFRadical.FullDefectRingEJZUnconditional.manuscriptOneSidedRingMaximalIsometryAllCharacteristics}
\leanverified{Manuscript/OneSidedMFRadical/FullDefectRingEJZUnconditional}{GroupApproximation.Manuscript.OneSidedMFRadical.FullDefectRingEJZUnconditional.manuscriptOneSidedRingMaximalReducedCStarRankTwoAllCharacteristics}
\leanverified{Manuscript/OneSidedMFRadical/FullDefectRingEJZUnconditional}{GroupApproximation.Manuscript.OneSidedMFRadical.FullDefectRingEJZUnconditional.manuscriptUnitGroupNotMFAllCharacteristics}
\end{corollary}

\begin{proof}
Choose $s,t\in R$ with $ts=1\ne st$, and let $S$ be the unital subring
they generate.  Let $u$ and $L=\EL_3(S)$ be as in
Lemma~\ref{lem:ring-compression-cell} for $S$; then $L$ has
property~\textup{(T)}~\cite[Theorem~1.1]{EJZ}, and the compression is
strict, since $e(sat)=0$ while $e\cdot1=e\ne0$ puts
$e_{12}(1)$ in $L\setminus uLu^{-1}$.  By
Proposition~\ref{prop:max-infinite} applied to $L\le\EL_n(R)$ and $u$,
$C^*_{\max}(\EL_n(R))$ contains a proper isometry.

If $R\ne0$ satisfies the hypothesis of
Theorem~\ref{thm:full-defect-ring}, then $G=\EL_n(R)$ is countable and
nontrivial, and
by Theorem~\ref{thm:full-defect-ring} it is not MF\@.  The algebra
$C^*_{\mathrm r}(G)$ is separable, stably finite because its canonical
trace is faithful, and not MF by Lemma~\ref{prop:mf-residual-calculus}.
Finally, the map $\jmath$ from the proof of
Theorem~\ref{thm:full-defect-ring} is injective, since
$T_i\jmath(A)S_j=A_{ij}$, and sends $\EL_4(R)$ into $R^\times$; that
group is not MF, and restricting the maps $V_n$ shows that MF passes to
subgroups.  So $R^\times$ is not MF\@.
\end{proof}

\subsection{The binary example}

Put $R=L_{\F_2}(1,2)$ and $H=R^\times$.  By~\eqref{eq:leavitt}, the
maps $x\mapsto(t_0x,t_1x)$ and $(y,z)\mapsto s_0y+s_1z$ are mutually
inverse isomorphisms of right $R$-modules between $R$ and $R\oplus R$,
and $H\cong\GL_4(R)=\EL_4(R)$~\cite[Proposition~4.2 and
Corollary~4.4]{KhanhThanh}.  We identify $H$ with $\EL_4(R)$.

\begin{proof}[Proof of Theorem~\ref{thm:headline}]
The relations give $t_0s_0=1$ and $t_1(1-s_0t_0)s_1=1$, so $R$
satisfies the hypothesis of Theorem~\ref{thm:full-defect-ring}, and
every homomorphism from $H$ to an MF group is trivial.  The ring $R$ is
finitely generated, so $H$ has property~\textup{(T)}
by~\cite[Theorem~1.1]{EJZ}, and hence is finitely
generated~\cite[Theorem~1.3.1]{BHV}.  Since $e_{12}(1)\ne1$, the group
$H$ is nontrivial.

Let $N$ be a normal subgroup of $H$.  The ring $R$ is purely infinite
simple~\cite{AbramsAranda} and therefore an exchange
ring~\cite{AraExchange}, so by Preusser's sandwich
theorem~\cite[Theorem~3]{Preusser} there is an ideal $I$ of $R$ with
\[
 \EL_4(R,I)\le N\le C_4(R,I),
\]
where $\EL_4(R,I)$ is the normal closure in $\EL_4(R)$ of the
elementary matrices $e_{ij}(a)$ with $a\in I$, and $C_4(R,I)$ is the
preimage of the center of $\GL_4(R/I)$.  Since $R$ is simple, either
$I=R$, and then $N\ge\EL_4(R)=H$, or $I=0$, and then every $g\in N$ is
central in $\GL_4(R)$.  In the second case, commuting with each
$e_{ij}(1)$ forces $g=\lambda I_4$ with $\lambda\in R^\times$, and
commuting with each $e_{ij}(a)$ then gives $\lambda a=a\lambda$ for all
$a\in R$.  So $\lambda$ lies in
$Z(R)=\F_2$~\cite[Corollary~4.3]{ArandaCrow}, and $N=1$.  So $H$ is
simple.  It is nontrivial, and every homomorphism from it to an MF
group is trivial, so $H$ is not MF\@.

Since $R\ne0$ satisfies the hypothesis of
Theorem~\ref{thm:full-defect-ring}, Corollary~\ref{cor:one-sided-ring-maximal}
shows that $C^*_{\mathrm r}(H)$ is separable, stably finite, and not
MF, and that $C^*_{\max}(H)$ contains a proper isometry.
\end{proof}

All rings in the following discussion are unital and associative.
Two idempotents $e,f$ of a ring are \emph{equivalent} if $e=xy$ and
$f=yx$ for some $x,y$; an idempotent $e$ is \emph{infinite} if
$e=f+g$ for orthogonal idempotents $f,g$ with $f$ equivalent to $e$ and
$g\ne0$; and a simple ring is \emph{purely infinite} if every nonzero
right ideal contains an infinite idempotent~\cite[Definitions~1.2]{AGP}.
The algebras $L_k(1,d)$ are purely infinite simple~\cite{AbramsAranda}.
Let $\GL(R)$ and $\EL(R)$ be the direct limits of $\GL_n(R)$ and
$\EL_n(R)$ along $A\mapsto\operatorname{diag}(A,1)$; the subgroup
$\EL(R)$ is normal in $\GL(R)$ by Whitehead's lemma, and
$K_1(R)=\GL(R)/\EL(R)$.  Write $\kappa\colon R^\times\to K_1(R)$ for the
canonical map.

\begin{theorem}[MF quotients of unit groups]\label{thm:mf-quotient-units}
Let $R$ be a countable unital purely infinite simple ring and let $n\ge1$.
Then every homomorphism from $\GL_n(R)$ to an MF group factors uniquely
through the canonical map $\GL_n(R)\to K_1(R)$, and $K_1(R)$, a
countable abelian group, is MF\@.  Equivalently,
$\Rad_{\mathrm{MF}}(\GL_n(R))=[\GL_n(R),\GL_n(R)]$ and
$\GL_n(R)/\Rad_{\mathrm{MF}}(\GL_n(R))\cong K_1(R)$.
\end{theorem}

\begin{proof}
The ring $M_n(R)$ is again countable, purely infinite, and
simple~\cite[Corollary~1.7]{AGP}, and $K_1(M_n(R))\cong K_1(R)$ by
Morita invariance, so it suffices to treat $n=1$.  Write $H=R^\times$
and $N=\Rad_{\mathrm{MF}}(H)$.  Ara, Goodearl, and Pardo show that $\kappa$ is surjective
with kernel $[H,H]$~\cite[Theorem~2.4]{AGP}.  A countable abelian group
$A$ is MF: $C^*_{\max}(A)$ is commutative and separable, hence
residually finite-dimensional and MF~\cite{BK}, and
Lemma~\ref{prop:mf-residual-calculus} applies.  So $\kappa$ is a homomorphism to an MF group, and $N\le\ker\kappa$.

For the converse we use Theorem~\ref{thm:full-defect-ring} in the
following form.  Every nonzero idempotent $e$ of $R$ is
infinite~\cite[Proposition~1.5]{AGP}, so $eRe$ contains $s,t$ with
$ts=e\ne st$, and $eRe$ is simple, so $e-st$ generates $eRe$ as a
two-sided ideal; hence every homomorphism from $\EL_m(eRe)$, $m\ge2$,
to an MF group is trivial.  Consequently, if $e$ and $P$ are nonzero
idempotents and $\theta\colon M_m(eRe)\to PRP$ is a ring isomorphism,
then
\begin{equation}\label{eq:corner-units}
 1-P+\theta(A)\in N\qquad(A\in\EL_m(eRe)),
\end{equation}
since these units form the image of $\EL_m(eRe)$ under a homomorphism
to $H$.

We use two facts from the proof of~\cite[Theorem~2.4]{AGP}, which
rest on~\cite[Proposition~1.5 and Corollary~1.7]{AGP} and on the
elementary reduction of Menal and Moncasi~\cite[proof of Theorem~2.2,
and the Remark after Corollary~2.3]{MenalMoncasi}.
\begin{enumerate}
\item[(a)] There are orthogonal idempotents $e_1,\dots,e_m$ with sum
$1$, $m\ge4$, such that $e_1,\dots,e_{m-1}$ are pairwise equivalent and
$e_m$ is equivalent to an idempotent $f\le e_1$.  Then $R$ is the ring
of $m\times m$ matrices over $T=e_1Re_1$ whose last column has entries
in $Tf$ and whose last row has entries in $fT$, and every unit $u$ of
$R$ factors as $u=PvQ$ with $P$ and $Q$ products of elementary matrices
$e_{ij}(x)$ of this matrix ring and $v=e_1+(1-e_1)v(1-e_1)$.
\item[(b)] For a nonzero idempotent $e$ and $n\ge2$ there are an
idempotent $f<e$ equivalent to $e$ and orthogonal idempotents
$r_2,\dots,r_n\le e-f$, each equivalent to $1$, such that, with
$r_1=1-e+f$ and $P=r_1+\dots+r_n$, some ring isomorphism
$\theta\colon M_n(R)\to PRP$ sends $\operatorname{diag}(v,1,\dots,1)$
to $(1-e)v(1-e)+f+r_2+\dots+r_n$ for every unit
$v=e+(1-e)v(1-e)$.
\end{enumerate}

Let $u\in\ker\kappa$ and write $u=PvQ$ as in (a).  The elementary
matrices with $i,j\le m-1$ form $\EL_{m-1}(T)$ inside
$(1-e_m)R(1-e_m)\cong M_{m-1}(T)$, so they lie in $N$
by~\eqref{eq:corner-units}; for distinct $i,j\le m-1$,
\[
 e_{im}(x)=[e_{ij}(1),e_{jm}(x)],\qquad
 e_{mj}(x)=[e_{mi}(x),e_{ij}(1)]
\]
are products of an element of $\EL_{m-1}(T)$ and a conjugate of its
inverse, so they lie in the normal subgroup $N$ as well.  Hence
$u\equiv v$ modulo $N$, and $\kappa(v)=\kappa(u)=0$ because
$N\le\ker\kappa$.

So let $v=e+(1-e)v(1-e)$ with $e\ne0$ and $\kappa(v)=0$.  By the
definition of $K_1$, $\operatorname{diag}(v,1,\dots,1)\in\EL_n(R)$ for
some $n\ge2$.  With $\theta$ and $P$ as in (b),
$v=1-P+\theta(\operatorname{diag}(v,1,\dots,1))$, which lies in $N$
by~\eqref{eq:corner-units}.  Hence $\ker\kappa\le N$.
\end{proof}

\begin{corollary}[Leavitt unit groups]\label{cor:leavitt-mf-quotient}
Let $k$ be a countable field, let $d\ge2$, let $R=L_k(1,d)$, and let
$H=R^\times$.  Then $H\cong\GL_d(R)$, $\Rad_{\mathrm{MF}}(H)=[H,H]=\EL_d(R)$, and
\[
 H/\EL_d(R)\cong K_1(R)\cong k^\times/(k^\times)^{d-1}.
\]
\end{corollary}

\begin{proof}
The maps $x\mapsto(t_1x,\dots,t_dx)$ and
$(y_1,\dots,y_d)\mapsto\sum_is_iy_i$ are mutually inverse isomorphisms
of right $R$-modules between $R$ and $R^d$, so $R\cong M_d(R)$ and
$H\cong\GL_d(R)$.  Theorem~\ref{thm:mf-quotient-units} identifies the
radical with $[H,H]$ and the quotient with $K_1(R)$.  Khanh--Thanh
show that $\GL_d(R)=\EL_d(R)D_d(k)$, where $D_d(k)$ is the abelian group
of diagonal matrices with entries in $k^\times$, which normalizes
$\EL_d(R)$, and that $K_1(R)\cong k^\times/(k^\times)^{d-1}$~\cite[proof
of Theorem~7.2]{KhanhThanh}.  So $[H,H]\le\EL_d(R)$, while
$\EL_d(R)\le[H,H]$ because every homomorphism from $\EL_d(R)$ to an MF
group is trivial (Theorem~\ref{thm:full-defect-ring}).
\end{proof}

For $d=2$ the quotient is trivial for every countable field $k$, so
every homomorphism from $L_k(1,2)^\times$ to an MF group is trivial.
For $k=\mathbb F_q$ the quotient is cyclic of order $\gcd(q-1,d-1)$.

\section{Amenable traces and Clifford lamps}
\label{sec:amenable-nonqd}

Brown observed that it was unknown whether every amenable trace on a
$C^*$-algebra is quasidiagonal~\cite[discussion preceding
Proposition~3.5.1]{Brown}.  For a separable unital $C^*$-algebra $A$ we
use the sequential form of his definitions.  A tracial state $\tau$ on
$A$ is \emph{amenable} if there are u.c.p. maps
$\phi_n\colon A\to M_{d_n}(\C)$ such that
\[
 \hsnorm{\phi_n(ab)-\phi_n(a)\phi_n(b)}\longrightarrow0,
 \qquad \tr_{d_n}(\phi_n(a))\longrightarrow\tau(a)
 \quad(a,b\in A).
\]
It is \emph{quasidiagonal} if the first limit holds in operator norm instead
of normalized Hilbert--Schmidt norm, with the same trace convergence.

For a countable group $G$, let $u_g$ denote the canonical unitary in
$C^*_{\max}(G)$.  The canonical trace $\tau_G$ is
determined by
\[
 \tau_G(u_g)=\begin{cases}1,&g=1,\\0,&g\ne1.\end{cases}
\]

\begin{theorem}[canonical-trace obstruction]
\label{thm:factorization-nonmf-trace}
Let $G$ be a countable group.  If $G$ is not MF, then its canonical trace
$\tau_G$ on $C^*_{\max}(G)$ is not quasidiagonal.  In particular, if
$\tau_G$ is amenable, then $\tau_G$ is an amenable trace that is not
quasidiagonal.
\leanverified{Manuscript/NinetyNineProblems/ProblemXGroups}{GroupApproximation.NinetyNineProblems.manuscriptCanonicalMaximalTraceNotIsQuasidiagonalTraceOfNotIsOperatorMF}
\leanverified{Manuscript/NinetyNineProblems/ProblemXGroups}{GroupApproximation.NinetyNineProblems.canonicalMaximalTrace_amenable_not_isQuasidiagonalTrace_of_not_isOperatorMF}
\end{theorem}

\begin{proof}
Suppose that $\tau_G$ is quasidiagonal, with u.c.p.\ maps
$\phi_n\colon C^*_{\max}(G)\to M_{d_n}(\C)$ as in the definition.
Asymptotic multiplicativity in operator norm makes
$a\mapsto[(\phi_n(a))_n]$ a unital $*$-homomorphism
$\Phi\colon C^*_{\max}(G)\to\mathcal Q_{\mathbf d}$.  If $\Phi(u_g)=1$,
then $\opnorm{\phi_n(u_g)-1}\to0$, so $\tr_{d_n}(\phi_n(u_g))\to1$,
while $\tr_{d_n}(\phi_n(u_g))\to\tau_G(u_g)$; so $g=1$.  Hence
$g\mapsto\Phi(u_g)$ is a corona homomorphism injective on $G$, and $G$
is MF, contrary to the hypothesis.
\end{proof}

A group $N$ is \emph{locally residually finite} if each of its finitely
generated subgroups is residually finite.

\begin{proposition}[amenable extensions of locally residually finite groups]
\label{prop:locally-rf-by-z-trace}
Let $1\to N\to G\to A\to1$ be an extension of countable groups with $N$
locally residually finite and $A$ amenable.  Then the canonical trace
of $C^*_{\max}(N)$ is quasidiagonal, and the canonical trace of
$C^*_{\max}(G)$ is amenable.
\leanverified{Analysis/AmenableExtensionAmenableTrace}{GroupApproximation.AmenableExtensionTrace.manuscriptPrintedAmenableExtensionTrace}
\end{proposition}

\begin{proof}
Write $g\mapsto\bar g$ for the quotient map
$G\to A$, and fix a finite set $\Sigma\subseteq G$, a finite F\o lner set
$F\subseteq A$, and a section $\sigma\colon A\to G$.  For $g\in \Sigma$ and
$x\in F$ put
\[
 b(g,x)=\sigma(\bar gx)^{-1}g\sigma(x)\in N .
\]
The subgroup $N_0$ generated by these finitely many elements is
residually finite; choose a finite quotient $\theta\colon N_0\to \Delta$
with $\theta(b(g,x))\ne1$ whenever $b(g,x)\ne1$, put $\Lambda=\ker\theta$,
and choose representatives $r_q\in N_0$ of the cosets of $\Lambda$ in $N_0$.
The cosets $\sigma(x)r_q\Lambda$, $x\in F$, $q\in \Delta$, are pairwise distinct,
since equality forces $x=x'$ in $A$ and then $q=q'$; let $T$ be their
set, let $\Pi$ be the projection of $\ell^2(G/\Lambda)$ onto $\ell^2(T)$, and
put $\Phi(b)=\Pi\lambda(b)\Pi|_{\ell^2(T)}$ for $b\in C^*_{\max}(G)$, where
$\lambda$ is the quasi-regular representation.  This map is u.c.p., and
for $g\in \Sigma$ and $\bar gx\in F$,
\[
 g\,\sigma(x)r_q\Lambda=\sigma(\bar gx)\,r_{\theta(b(g,x))q}\Lambda\in T .
\]
So for $g,h\in \Sigma$ the identity
$\Phi(u_{gh})-\Phi(u_g)\Phi(u_h)=\Pi\lambda(g)(1-\Pi)\lambda(h)\Pi$ and the
rank of $(1-\Pi)\lambda(h)\Pi$, at most $|\{x\in F:\bar hx\notin F\}|\,|\Delta|$,
give
\[
 \hsnorm{\Phi(u_{gh})-\Phi(u_g)\Phi(u_h)}^2
 \le\frac{|\{x\in F:\bar hx\notin F\}|}{|F|}.
\]
The normalized trace of $\Phi(u_g)$ is the fraction of points of $T$
fixed by $g$, and for $g\in \Sigma\setminus\{1\}$ it is zero: if
$\bar g\ne1$, then $g$ moves the $F$-coordinate of every point, and if
$\bar g=1$, then $g\sigma(x)r_q\Lambda=\sigma(x)r_{\theta(b(g,x))q}\Lambda$ with
$b(g,x)=\sigma(x)^{-1}g\sigma(x)\ne1$, so $\theta(b(g,x))q\ne q$.
Taking F\o lner sets with $|\{x\in F:\bar hx\notin F\}|/|F|\to0$ for
the finitely many $h\in \Sigma$, along an exhaustion of $G$ by finite sets
$\Sigma$, the maps $\Phi$ are asymptotically multiplicative in normalized
Hilbert--Schmidt norm and their traces converge to $\tau_G$ on the
canonical unitaries; density and contractivity extend both limits to
$C^*_{\max}(G)$.

For the canonical trace of $C^*_{\max}(N)$, run the same construction
with $N$ in place of $G$, $A=1$, and $F=\{1\}$: no boundary term
appears, so the maps $\Phi$ are multiplicative on $\Sigma$ in operator norm,
and their traces vanish on $\Sigma\setminus\{1\}$; the same limits give
quasidiagonality.
\end{proof}

We now construct a group to which
Theorem~\ref{thm:compression-criterion} applies through a finite
central subgroup: its intrinsic defect subgroup contains a commutator
whose square is a central involution.  Let $\Gamma$ be a
countable group with property~\textup{(T)}, let
$\alpha\colon\Gamma\to\Gamma$ be injective but not surjective, and choose
$a\in\Gamma\setminus\alpha(\Gamma)$.  Put
\[
 T_\alpha=\varinjlim
   (\Gamma\xrightarrow{\alpha}\Gamma
    \xrightarrow{\alpha}\cdots),
 \qquad V=T_\alpha\rtimes\langle t\rangle,
\]
where $tgt^{-1}=\alpha(g)$ on the level-zero copy of $\Gamma$, so that
$V$ is the ascending HNN extension of $\Gamma$ along $\alpha$ and
$T_\alpha=\bigcup_{n\ge0}t^{-n}\Gamma t^n$.  Put
$X=V/\Gamma$, the left-coset space.  The Clifford lamp group
$\operatorname{Cl}(X)$ is the group with generators $\varepsilon$ and
$c_x$, $x\in X$, and relations
\[
 \varepsilon^2=c_x^2=1,\qquad [\varepsilon,c_x]=1,\qquad
 [c_x,c_y]=\varepsilon\quad(x\ne y).
\]
To see that $\varepsilon\ne1$, fix a total order on $X$, let $\Sigma$ be
the $\F_2$-vector space of finitely supported functions $X\to\F_2$, and
for $f,g\in \Sigma$ put $B(f,g)=\sum_{x>y}f(x)g(y)$.  Since $B$ is bilinear,
\[
 (a,f)(b,g)=\bigl(a+b+B(f,g),\ f+g\bigr)
\]
is a group law on $\F_2\times \Sigma$, in which $(1,0)$ is a central
involution, each $(0,\delta_x)$ is an involution, and for $x\ne y$
exactly one of $B(\delta_x,\delta_y)$ and $B(\delta_y,\delta_x)$ equals
$1$, so the commutator of $(0,\delta_x)$ and $(0,\delta_y)$ is $(1,0)$.
So $\varepsilon\mapsto(1,0)$, $c_x\mapsto(0,\delta_x)$ defines a
homomorphism $\operatorname{Cl}(X)\to\F_2\times \Sigma$.  The relations let
us write every element of $\operatorname{Cl}(X)$ as
\[
 \varepsilon^a c_{x_1}\cdots c_{x_r},
 \qquad a\in\F_2,\quad x_1<\cdots<x_r,
\]
and this word maps to $(a,\delta_{x_1}+\cdots+\delta_{x_r})$, so the
expression is unique and the homomorphism is an isomorphism.  In
particular, for a finite subset $Y\subseteq X$, the subgroup generated
by $\varepsilon$ and the $c_y$ with $y\in Y$ has order $2^{|Y|+1}$, so
$\operatorname{Cl}(X)$ is countable and locally finite.  Killing
$\varepsilon$ leaves $\Sigma\cong\bigoplus_XC_2$, so
$\operatorname{Cl}(X)$ is a central extension of $\bigoplus_XC_2$ by
$C_2$, and $W$ below is a central extension of the permutational wreath
product $C_2\wr_XV$ by $\langle\varepsilon\rangle$.  The relations
are invariant under permutations of $X$, so every permutation of $X$
induces an automorphism of $\operatorname{Cl}(X)$ that permutes the
$c_x$ accordingly and fixes $\varepsilon$.  In this way $V$ acts on
$\operatorname{Cl}(X)$ through its action on $X$, and we put
\begin{equation}\label{eq:affine-clifford-witness}
 W=\operatorname{Cl}(X)\rtimes V.
\end{equation}

\begin{proposition}[Clifford obstruction from a proper self-embedding]
\label{prop:clifford-self-embedding}
The group $W$ is not MF.  More precisely, every homomorphism from $W$ to an
MF group kills $\varepsilon$.
\leanverified{Sofic/CliffordWitnessDirectDefect}{GroupApproximation.CliffordWitnessDirectDefect.manuscriptCliffordWitnessNotIsOperatorMF}
\leanverified{Sofic/CliffordWitnessDirectDefect}{GroupApproximation.CliffordWitnessDirectDefect.manuscriptMapSignEqOneOfIsOperatorMFTarget}
\end{proposition}

\begin{proof}
Identify $\Gamma$ with its level-zero copy in $V\le W$, and let $c$ be
the lamp at the root coset $\Gamma\in X$; it centralizes $\Gamma$, which
fixes that coset.  Then $tct^{-1}$ is the lamp at $t\Gamma$ and
$a(tct^{-1})a^{-1}$ is the lamp at $at\Gamma$.  These cosets are
distinct, because $t\Gamma=at\Gamma$ would mean
$a\in t\Gamma t^{-1}=\alpha(\Gamma)$; so $x=tct^{-1}$ and $y=axa^{-1}$
are distinct lamps, involutions with $[x,y]=\varepsilon\ne1$.  Then
$d=[x,a]=xy$ and $d^2=xyxy=[x,y]=\varepsilon$, and $d=[tct^{-1},a]$ is a
generator of $\mathfrak D_W(\Gamma)$, since
$t\Gamma t^{-1}\le\Gamma$, $c\in C_W(\Gamma)$, and $a\in\Gamma$.  So $\langle\varepsilon\rangle\le\mathfrak D_W(\Gamma)$.
The subgroup $\langle\varepsilon\rangle$ is central, so normal, and it
is finite, so it has property~\textup{(T)}.  So
Theorem~\ref{thm:compression-criterion} applies to the countable group
$W$ with $L=\Gamma$ and $K=\langle\varepsilon\rangle$: every
homomorphism from $W$ to an MF group is trivial on $\varepsilon$.  In
particular, $W$ is not MF\@.
\end{proof}

\begin{proposition}[the locally residually finite case]
\label{prop:clifford-locally-rf}
Suppose in addition that $\Gamma$ is residually finite and that
$[\Gamma:\alpha(\Gamma)]<\infty$.  Then
\[
 W\cong W_0\rtimes\mathbb Z,
 \qquad W_0=\operatorname{Cl}(X)\rtimes T_\alpha,
\]
where $W_0$ is locally residually finite.  Consequently $W$ is sofic but not
MF, the canonical trace on $C^*_{\max}(W_0)$ is quasidiagonal, and the
canonical trace on $C^*_{\max}(W)$ is amenable and not quasidiagonal.
\leanverified{Sofic/CliffordWitnessSoficPrinted}{GroupApproximation.AmenableTraceTheorem.manuscriptCliffordLocallyRF}
\end{proposition}

\begin{proof}
Reassociating the semidirect products yields
\[
 W\cong W_0\rtimes\mathbb Z,
 \qquad W_0=\operatorname{Cl}(X)\rtimes T_\alpha.
\]
\leanverified{Sofic/CliffordWitnessLocallyRFByInt}{GroupApproximation.CliffordWitnessLocallyRFByInt.manuscriptAmbientEquivShiftKernelForByInt}
For $n\ge0$, write $\Gamma_n=t^{-n}\Gamma t^n$ for the image of the $n$th
copy of $\Gamma$ in $T_\alpha$.  Since $[\Gamma:\alpha(\Gamma)]<\infty$,
each $\Gamma_n$ has finite index in $\Gamma_{n+1}$; every element of
$T_\alpha$ lies in some $\Gamma_m$, and conjugation by $t$ maps
$\Gamma_{m+1}$ onto $\Gamma_m$, so $\Gamma$ is commensurated by $V$.
Hence the stabilizer in $\Gamma_n$ of every point of $X=V/\Gamma$ has
finite index, and every $\Gamma_n$-orbit in $X$ is finite.

A finite subset of $W_0$ involves only finitely many lamps $c_x$, and its
$T_\alpha$-coordinates lie in one $\Gamma_n$.  Let $Y$ be the union of
the $\Gamma_n$-orbits of these finitely many $x$.  Then $Y$ is finite
and $\Gamma_n$-invariant, and the finite subset is contained in
\[
 C_Y\rtimes\Gamma_n,
\]
where $C_Y=\langle\varepsilon,c_y:y\in Y\rangle$ is finite by the
normal form above.  This semidirect product is
residually finite.  Indeed, an element with nontrivial $\Gamma_n$-component
survives in a finite quotient of $\Gamma_n$.  For a nontrivial element of
$C_Y$, let $J$ be the kernel of the action
$\Gamma_n\to\operatorname{Aut}(C_Y)$.  The quotient $\Gamma_n/J$ is finite,
and the element has nontrivial image in the finite group
$C_Y\rtimes(\Gamma_n/J)$.  Therefore every finitely generated subgroup of $W_0$ is residually
finite, hence sofic.
\leanverified{Sofic/CliffordWitnessLocallyRFByInt}{GroupApproximation.CliffordWitnessLocallyRFByInt.manuscriptShiftKernelForIsLocallyResiduallyFinite}

The group $W_0$ is their directed union, and $W/W_0\cong\mathbb Z$
is amenable.  Soficity passes to directed unions and to extensions with
amenable quotient~\cite[Theorem~1]{ElekSzabo}, so $W$ is sofic.
Proposition~\ref{prop:clifford-self-embedding} shows that it is not MF.
By Proposition~\ref{prop:locally-rf-by-z-trace}, applied to the
extension $1\to W_0\to W\to\mathbb Z\to1$, the canonical trace of
$C^*_{\max}(W_0)$ is quasidiagonal and that of $C^*_{\max}(W)$ is
amenable; by Theorem~\ref{thm:factorization-nonmf-trace} the latter is
not quasidiagonal.
\end{proof}

For a concrete instance, take
\[
 \bar\Gamma=\mathbb Z^3\rtimes\mathrm{SL}_3(\mathbb Z)
 =\left\{
   \begin{pmatrix}A&v\\0&1\end{pmatrix}:
   A\in\mathrm{SL}_3(\mathbb Z),\ v\in\mathbb Z^3
  \right\}
 \le\GL_4(\mathbb Z).
\]
It is residually finite, since reduction modulo a suitable integer
separates any two distinct integral matrices, and it has
property~\textup{(T)}~\cite[Example~1.7.4(i)]{BHV}.
\leanverified{Sofic/CommutingLampCollapse}{GroupApproximation.CommutingLampCollapse.gammaBar_hasKazhdanPropertyT}
\leanverified{Monsters/ExplicitIntegralLinearModel}{GroupApproximation.ExplicitIntegralLinearModel.gammaBar_residuallyFinite}
Let
\[
 D=\operatorname{diag}(2,2,2,1),\qquad
 \alpha(g)=DgD^{-1},
\]
so that $\alpha(v,A)=(2v,A)$.  This is injective, and its image consists
of the affine matrices whose translation coordinates are all even, so
$[\bar\Gamma:\alpha(\bar\Gamma)]=8$.  Translation $a$ by the first
standard basis vector lies outside $\alpha(\bar\Gamma)$.  So the preceding
construction applies to $\bar\Gamma$, $\alpha$ and $a$.

\begin{proof}[Proof of Theorem~\ref{thm:amenable-trace}]
Apply Proposition~\ref{prop:clifford-locally-rf} to $\bar\Gamma$,
$\alpha$, and $a$.
\leanverified{Sofic/CliffordWitnessLocallyRFByInt}{GroupApproximation.CliffordWitnessLocallyRFByInt.witnessCanonicalMaximalTrace_isAmenableTrace}
\leanverified{Manuscript/NinetyNineProblems/ProblemXCliffordWitness}{GroupApproximation.NinetyNineProblems.witnessCanonicalTrace_amenable_not_quasidiagonal}
\end{proof}

\section{A torsion-free finitely presented example}
\label{sec:torsion-free}

For an acylindrically hyperbolic group $G$, choose the generating set $A$
provided by Hull~\cite[Theorem~3.12]{Hull}.  A subgroup is \emph{suitable}
if it acts non-elementarily on the Cayley graph $\operatorname{Cay}(G,A)$ and normalizes no nontrivial
finite subgroup~\cite[Definition~1.4]{Hull}.  Hull states his small cancellation
theorem~\cite[Theorem~7.1(a), (b), (c), and (e)]{Hull} with injectivity on
a ball of $\operatorname{Cay}(G,A)$; since every finite subset
of $G$ lies in some ball, the theorem takes the following form.

\begin{theorem}[Hull's small cancellation theorem]\label{thm:hull}
Let $G$ be acylindrically hyperbolic, let $N\le G$ be suitable with respect
to $A$, and let $g_1,\dots,g_m\in G$ and a finite subset $\Omega\subseteq G$
be given.  Then there is a surjective homomorphism $\varphi\colon G\to Q$
such that $Q$ is acylindrically hyperbolic, $\varphi|_\Omega$ is injective,
$\varphi(g_i)\in\varphi(N)$ for all $i$, and every element of finite order
in $Q$ is the image of an element of the same order in $G$.
\end{theorem}

Hull's proof treats $m=1$ by passing to $G/\normal{r}_G$ for one
element $r$ and the general case by induction on $m$~\cite[proof of Theorem~7.1]{Hull},
so $\ker\varphi$ is the normal closure of $m$ elements and $Q$ is finitely
presented when $G$ is.

\begin{lemma}[saturation]\label{lem:saturation}
Let $G$ be finitely presented, torsion-free, and acylindrically hyperbolic,
let $N\trianglelefteq G$ be nontrivial, and let $\Omega\subseteq G$ be finite.
Then there is a surjective homomorphism $\varphi\colon G\to Q$ such that $Q$
is two-generated, finitely presented, torsion-free, and acylindrically
hyperbolic, $\varphi|_\Omega$ is injective, and $\varphi(N)=Q$.
\end{lemma}

\begin{proof}
The subgroup $N$ is normal, and it is infinite because $G$ is
torsion-free and $N\ne1$, so it acts non-elementarily on
$\operatorname{Cay}(G,A)$ by Osin~\cite[Lemma~7.1]{Osin}; since $G$ is torsion-free,
$N$ normalizes no nontrivial finite subgroup, so $N$ is suitable.  By Hull~\cite[Corollary~5.7 and Lemma~5.8]{Hull}, $N$ contains
two elements $h_1,h_2$ such that $N_0=\langle h_1,h_2\rangle$ is again
suitable, possibly with respect to a larger generating set
$A'\supseteq A$, and Theorem~\ref{thm:hull} holds for $A'$ as well,
since every finite subset of $G$ lies in a ball of $\operatorname{Cay}(G,A')$.  Apply
Theorem~\ref{thm:hull} to $N_0$ with $g_1,\dots,g_m$ a
finite generating set of $G$ and with the set $\Omega$.  Then
$Q=\langle\varphi(g_1),\dots,\varphi(g_m)\rangle\le\varphi(N_0)\le\varphi(N)\le Q$,
so $Q=\varphi(N_0)=\langle\varphi(h_1),\varphi(h_2)\rangle$ and
$\varphi(N)=Q$.  The quotient is torsion-free by Theorem~\ref{thm:hull} and
finitely presented because $\ker\varphi$ is normally generated by $m$
elements.
\end{proof}

Fournier-Facio constructs a finitely presented torsion-free group $G_0$
with property~\textup{(T)}, a subgroup $\Gamma\le G_0$ with
property~\textup{(T)}, an element $t\in G_0$ with
$t\Gamma t^{-1}\le\Gamma$, and a subgroup $J\le G_0$ isomorphic to a
finitely presented infinite simple group, such that $[\Gamma,J]=1$ and
$tJt^{-1}\le\Gamma$~\cite[\S2]{FFF}.  The group $G_0$ is obtained there
as a common quotient of two finitely generated acylindrically hyperbolic
groups by Hull's theorem~\cite[Corollary~7.4]{Hull}, which allows the
quotient to be chosen acylindrically hyperbolic; we take $G_0$ to be
such a quotient.  Put $S=tJt^{-1}$.  Since $J$ is simple and
nonabelian, $J$ and $S$ are perfect.

For $c\in J$ and $\ell\in S$, the commutator $[tct^{-1},\ell]$ lies in
$\mathfrak D_{G_0}(\Gamma)$ by~\eqref{eq:intrinsic-defect}, since $t$
conjugates $\Gamma$ into itself, $c$ centralizes $\Gamma$, and
$S\le\Gamma$.  As $c$ ranges over $J$, the element $tct^{-1}$ ranges
over $S$, so $[S,S]\le\mathfrak D_{G_0}(\Gamma)$, and $S$ is perfect, so
\[
 S\le\mathfrak D_{G_0}(\Gamma).
\]

\begin{proof}[Proof of Theorem~\ref{thm:torsion-free}]
Let $N$ be the normal closure of $S$ in $G_0$; it is nontrivial because
$S$ is.  By Lemma~\ref{lem:saturation} applied to $G_0$, $N$, and
$\Omega=\varnothing$, there is a surjective homomorphism
$\varphi\colon G_0\to Q$ with $Q$ two-generated, finitely presented,
torsion-free, and acylindrically hyperbolic, and $\varphi(N)=Q$.  The
group $Q$ is infinite because it is acylindrically hyperbolic, and it has
property~\textup{(T)} as a quotient of $G_0$.

By~\eqref{eq:defect-functorial}, $\varphi(S)\le\mathfrak D_Q(\varphi(\Gamma))$.  The subgroup
$\mathfrak D_Q(\varphi(\Gamma))$ is normal in $Q$, and the normal
closure of $\varphi(S)$ is $\varphi(N)=Q$, so
\[
 \mathfrak D_Q(\varphi(\Gamma))=Q.
\]
Both $Q$ and $\varphi(\Gamma)$ have property~\textup{(T)}, as quotients of
$G_0$ and $\Gamma$, respectively.  By the last assertion of
Theorem~\ref{thm:compression-criterion}, every homomorphism from $Q$ to
an MF group is trivial.  If a quotient $\bar Q$ of $Q$ is MF, then the
quotient map $Q\to\bar Q$ is trivial, so $\bar Q=1$.
\end{proof}

\begin{corollary}\label{cor:regular-nonmf-algebra}
The algebra $C^*_{\mathrm r}(Q)$ is separable, unital, generated by two
unitaries, simple, has a unique tracial state, has stable rank one, is
stably finite, and is not MF\@.
\end{corollary}

\begin{proof}
The group $Q$ is countable, torsion-free, and acylindrically
hyperbolic, so it contains a non-degenerate hyperbolically embedded
subgroup~\cite[Theorem~1.2]{Osin} and has no nontrivial finite normal
subgroup.  Dahmani, Guirardel, and Osin give simplicity and uniqueness
of the trace~\cite[Theorem~2.35]{DGO}, and Gerasimova and Osin give
density of the invertible elements, which is stable rank
one~\cite[Theorem~1.1]{GO}.  The algebra is separable and generated by
the canonical unitaries of the two generators; its canonical trace is
faithful, so it is stably finite; and it is not MF by
Lemma~\ref{prop:mf-residual-calculus}, since $Q$ is not MF
(Theorem~\ref{thm:torsion-free}).
\end{proof}

\section*{Origin and authorship}

In early August 2026 the author was using AI models on two questions
raised by the nonsofic group of~\cite{OAI}: whether hyperlinearity
implies soficity, and whether every group is hyperlinear.  In one ChatGPT conversation, GPT-5.6 Sol, on ``high''
thinking mode, worked on weak MF (matricial field) approximations of a
Clifford-twisted wreath product as a step toward hyperlinearity.  Earlier in that
conversation the author had asked whether to pivot to the MF question,
and Sol had advised against it.  The author gave Sol the existing
progress on the hyperlinear problem and sent this prompt:
\begin{quote}
``Achieve f\{MF of the Clifford Kun--Thom extension \}$\widetilde W$
in such a way that it archives our goals and ends everything, or,
possibly, a different breakthrough, leading to achieving the final
result.  You got this!''
\end{quote}
Sol answered with an MF counterexample: every norm-matrix model of the
Clifford-twisted wreath product kills the Clifford sign.  That was on
August~9.  Also on August~9, the models proved that the maximal group
$C^*$-algebra of a strictly compressed Kazhdan group is infinite.
Parallel chats, including one in ``Pro'' mode, and the agents working
on the hyperlinear question, did not find the counterexample.  By
August~12 the models had reduced the obstruction to one commutator in
a finitely presented group and had proved in Lean that it is not MF\@.  They wrote the
first draft of this paper that evening.

The author prompted the models, set the direction, retrieved the
literature, and suggested the MF question.  With the models, the
author explored several routes to a non-MF group, shaped the current
proof, and checked and revised the citations.  The models drafted much of the text;
the author wrote the rest and edited the whole for clarity and for the
mathematical presentation.  GPT-5.6 Sol in Codex, and Claude
Fable~5 and Claude Opus~5 in Claude Code, wrote the Lean proofs.  The
whole exchange is in the repository's public commit
history.\footnote{\url{https://github.com/SauersML/group-approximation}}

The author consulted experts from August~13, and their comments shaped
the revisions.  Many models also reviewed drafts; Astra's reviews
shortened the paper and simplified the proof of
Theorem~\ref{thm:normal-kazhdan}.  We recommend Eckhardt's
writeup~\cite{Eckhardt}: a short, elementary proof that certain
generalized wreath products are not MF\@.

\section*{Acknowledgments}

The author thanks Francesco Fournier-Facio for early advice and detailed
comments on earlier drafts, and Alon Dogon for organizing the ensuing
discussion and for comments on the surrounding literature.  The author also
thanks Tatiana Shulman for clarifying the terminology and historical context
of the MF question, Caleb Eckhardt for discussions of MF representations
and matrix norms and for generously writing and sharing the expository
note~\cite{Eckhardt}, Rufus Willett for reading the arguments, and Bruce
Blackadar for correspondence on the origin of the MF and NF problems.

\begin{thebibliography}{99}
\bibitem[AA]{AbramsAranda}
G.~Abrams and G.~Aranda Pino,
\emph{Purely infinite simple Leavitt path algebras},
J. Pure Appl. Algebra \textbf{207} (2006), no.~3, 553--563,
\doi{10.1016/j.jpaa.2005.10.010}.

\bibitem[AGP]{AGP}
P.~Ara, K.~R. Goodearl, and E.~Pardo,
\emph{$K_0$ of purely infinite simple regular rings},
$K$-Theory \textbf{26} (2002), 69--100,
\href{https://arxiv.org/abs/math/0111066}{arXiv:math/0111066}.

\bibitem[AW]{AkemannWalter}
C.~A. Akemann and M.~E. Walter,
\emph{Unbounded negative definite functions},
Canad. J. Math. \textbf{33} (1981), no.~4, 862--871.

\bibitem[AC]{ArandaCrow}
G.~Aranda Pino and K.~Crow,
\emph{The center of a Leavitt path algebra},
Rev. Mat. Iberoam. \textbf{27} (2011), no.~2, 621--644,
\doi{10.4171/RMI/649}.

\bibitem[Ara]{AraExchange}
P.~Ara,
\emph{The exchange property for purely infinite simple rings},
Proc. Amer. Math. Soc. \textbf{132} (2004), no.~9, 2543--2547,
\doi{10.1090/S0002-9939-04-07369-1}.

\bibitem[BDL]{BDL}
B.~Bachner, A.~Dogon, and A.~Lubotzky,
\emph{On $L^1$-approximation of groups},
J. Algebra \textbf{702} (2026), 235--243,
\doi{10.1016/j.jalgebra.2026.04.041},
\href{https://arxiv.org/abs/2508.17392}{arXiv:2508.17392}.

\bibitem[BHV]{BHV}
B.~Bekka, P.~de~la~Harpe, and A.~Valette,
\emph{Kazhdan's property \textup{(T)}},
New Mathematical Monographs, vol.~11,
Cambridge University Press, Cambridge, 2008,
\doi{10.1017/CBO9780511542749}.

\bibitem[BK]{BK}
B.~Blackadar and E.~Kirchberg,
\emph{Generalized inductive limits of finite-dimensional $C^*$-algebras},
Math. Ann. \textbf{307} (1997), no.~3, 343--380,
\doi{10.1007/s002080050039}.

\bibitem[BK94]{BK94}
B.~Blackadar and E.~Kirchberg,
\emph{Generalized inductive limits of finite-dimensional $C^*$-algebras},
in \emph{$C^*$-Algebren}, Tagungsbericht 6/1994,
Mathematisches Forschungsinstitut Oberwolfach, 1994, p.~2.

\bibitem[Br]{Brown}
N.~P. Brown,
\emph{Invariant means and finite representation theory of $C^*$-algebras},
Mem. Amer. Math. Soc. \textbf{184} (2006), no.~865, viii+105,
\doi{10.1090/memo/0865}.

\bibitem[CDE]{CDE}
J.~R. Carri\'on, M.~Dadarlat, and C.~Eckhardt,
\emph{On groups with quasidiagonal $C^*$-algebras},
J. Funct. Anal. \textbf{265} (2013), no.~1, 135--152,
\doi{10.1016/j.jfa.2013.04.004}.

\bibitem[DGO]{DGO}
F.~Dahmani, V.~Guirardel, and D.~Osin,
\emph{Hyperbolically embedded subgroups and rotating families in groups
acting on hyperbolic spaces},
Mem. Amer. Math. Soc. \textbf{245} (2017), no.~1156,
\doi{10.1090/memo/1156}.

\bibitem[Ec]{Eckhardt}
C.~Eckhardt,
\emph{Non-MF groups and non-finite full group $C^*$-algebras},
preprint, 2026,
\href{https://arxiv.org/abs/2608.28772}{arXiv:2608.28772}.

\bibitem[EJZ]{EJZ}
M.~Ershov and A.~Jaikin-Zapirain,
\emph{Property~\textup{(T)} for noncommutative universal lattices},
Invent. Math. \textbf{179} (2010), 303--347,
\doi{10.1007/s00222-009-0218-2}.

\bibitem[ES]{ElekSzabo}
G.~Elek and E.~Szab\'o,
\emph{On sofic groups},
J. Group Theory \textbf{9} (2006), no.~2, 161--171,
\doi{10.1515/JGT.2006.011},
\href{https://arxiv.org/abs/math/0305352}{arXiv:math/0305352}.

\bibitem[FF]{FFF}
F.~Fournier-Facio,
\emph{A torsion-free non-sofic group},
preprint, 2026,
\href{https://arxiv.org/abs/2608.02025}{arXiv:2608.02025}.

\bibitem[GKM]{GaoEtAl}
D.~Gao, S.~Kunnawalkam Elayavalli, A.~Manzoor, and G.~Patchell,
\emph{A new source of purely finite matricial fields},
preprint, 2026,
\href{https://arxiv.org/abs/2603.24502}{arXiv:2603.24502}.

\bibitem[GO]{GO}
M.~Gerasimova and D.~Osin,
\emph{On invertible elements in reduced $C^*$-algebras of acylindrically
hyperbolic groups},
J. Funct. Anal. \textbf{279} (2020), no.~7, 108689,
\doi{10.1016/j.jfa.2020.108689},
\href{https://arxiv.org/abs/1910.14524}{arXiv:1910.14524}.

\bibitem[GH]{GoldbringHart}
I.~Goldbring and B.~Hart,
\emph{The universal theory of the hyperfinite $\mathrm{II}_1$ factor is not
computable},
Bull. Symbolic Logic \textbf{30} (2024), no.~2, 181--198,
\doi{10.1017/bsl.2024.7},
\href{https://arxiv.org/abs/2006.05629}{arXiv:2006.05629}.

\bibitem[Hu]{Hull}
M.~Hull,
\emph{Small cancellation in acylindrically hyperbolic groups},
Groups Geom. Dyn. \textbf{10} (2016), no.~4, 1077--1119,
\doi{10.4171/GGD/377}.

\bibitem[JNVWY]{MIPRE}
Z.~Ji, A.~Natarajan, T.~Vidick, J.~Wright, and H.~Yuen,
\emph{$\operatorname{MIP}^*=\operatorname{RE}$},
Commun. ACM \textbf{64} (2021), no.~11, 131--138,
\doi{10.1145/3485628},
\href{https://arxiv.org/abs/2001.04383}{arXiv:2001.04383}.

\bibitem[K--T]{KhanhThanh}
H.~V. Khanh and V.~H. Thanh,
\emph{Matrix generators for the unit groups of $L_K(1,d)$},
preprint, 2026,
\href{https://arxiv.org/abs/2607.10351}{arXiv:2607.10351}.

\bibitem[Ku16]{Kun16}
G.~Kun,
\emph{On sofic approximations of property~\textup{(T)} groups},
preprint, 2016,
\href{https://arxiv.org/abs/1606.04471}{arXiv:1606.04471}.

\bibitem[KT19]{KT19}
G.~Kun and A.~Thom,
\emph{Inapproximability of actions and Kazhdan's property~\textup{(T)}},
preprint, 2019,
\href{https://arxiv.org/abs/1901.03963}{arXiv:1901.03963}.

\bibitem[Kor]{Korchagin}
A.~I. Korchagin,
\emph{The MF-property for countable discrete groups},
Math. Notes \textbf{102} (2017), no.~2, 198--211,
\doi{10.1134/S0001434617070227},
\href{https://arxiv.org/abs/1704.06906}{arXiv:1704.06906}.

\bibitem[Lev]{Leavitt}
W.~G.~Leavitt,
\emph{The module type of a ring},
Trans. Amer. Math. Soc. \textbf{103} (1962), 113--130,
\doi{10.1090/S0002-9947-1962-0132764-X}.

\bibitem[LO]{LubotzkyOppenheim}
A.~Lubotzky and I.~Oppenheim,
\emph{Non $p$-norm approximated groups},
J. Anal. Math. \textbf{141} (2020), no.~1, 305--321,
\doi{10.1007/s11854-020-0119-2}.

\bibitem[MM]{MenalMoncasi}
P.~Menal and J.~Moncasi,
\emph{$K_1$ of von Neumann regular rings},
J. Pure Appl. Algebra \textbf{33} (1984), no.~3, 295--312,
\doi{10.1016/0022-4049(84)90064-1}.

\bibitem[OAI]{OAI}
OpenAI,
\emph{Nonsofic groups exist},
in \emph{Ten advances in mathematics and theoretical computer science},
Chapter~3, 78--95, 2026,
\url{https://cdn.openai.com/pdf/ten-proofs-oai.pdf}.

\bibitem[Os]{Osin}
D.~Osin,
\emph{Acylindrically hyperbolic groups},
Trans. Amer. Math. Soc. \textbf{368} (2016), no.~2, 851--888,
\doi{10.1090/tran/6343}.

\bibitem[Pre]{Preusser}
R.~Preusser,
\emph{On general linear groups over exchange rings},
Linear and Multilinear Algebra \textbf{70} (2022), no.~4, 705--713,
\doi{10.1080/03081087.2020.1743636},
\href{https://arxiv.org/abs/1912.11386}{arXiv:1912.11386}.

\bibitem[Sch]{Schafhauser}
C.~Schafhauser,
\emph{Finite dimensional approximations of certain amalgamated free
products of groups},
Groups Geom. Dyn. \textbf{20} (2026), no.~2, 607--615,
\doi{10.4171/GGD/826},
\href{https://arxiv.org/abs/2306.02498}{arXiv:2306.02498}.

\bibitem[STW]{STW}
C.~Schafhauser, A.~Tikuisis, and S.~White,
\emph{Nuclear $C^*$-algebras: 99 problems},
to appear in M\"unster J. Math.,
\href{https://arxiv.org/abs/2506.10902}{arXiv:2506.10902}.

\bibitem[Th]{Thom}
A.~Thom,
\emph{Finitary approximations of groups and their applications},
Proceedings of the International Congress of Mathematicians---Rio de
Janeiro 2018,
\href{https://arxiv.org/abs/1712.01052}{arXiv:1712.01052}.

\bibitem[TWW]{TWW}
A.~Tikuisis, S.~White, and W.~Winter,
\emph{Quasidiagonality of nuclear $C^*$-algebras},
Ann. of Math. (2) \textbf{185} (2017), no.~1, 229--284,
\doi{10.4007/annals.2017.185.1.4}.

\end{thebibliography}
\end{document}
